% Revisão Sistemática da Literatura — PRISMA 2020
% Alvo primário: venue brasileiro (português, ABNT)
% Estrutura flexível para adaptação a venues internacionais
\documentclass[12pt,a4paper]{article}

% =====================================================================
% Encoding e idioma
% =====================================================================
\usepackage[utf8]{inputenc}
\usepackage[T1]{fontenc}
\usepackage[brazilian]{babel}
\usepackage{times}

% =====================================================================
% Tipografia refinada
% =====================================================================
% microtype: protrusion e expansion de caracteres — melhora line breaks,
% reduz overfull hboxes e dá acabamento profissional ao texto.
% Invisível mas é a maior diferença entre LaTeX amador e polido.
\usepackage[protrusion=true, expansion=true, tracking=false,
            disable=false]{microtype}
% Suprime patch de \footnote que entra em conflito com \thanks em
% \maketitle (warning benigno, mas ruidoso no log)
\microtypesetup{patch=none}

% Remover \sloppy: microtype resolve a maioria dos overfull hboxes sem
% produzir os espaços inter-palavra excessivos que \sloppy causa.
% Se um hbox persistir, ajustar a frase localmente.

% =====================================================================
% Cores — paleta consistente para todo o manuscrito
% =====================================================================
% xcolor com opção table para \rowcolor em tabelas
\usepackage[table]{xcolor}

% Paleta categórica: ColorBrewer Dark2 (colorblind-safe, 8 classes)
% Usada para tipos de arquitetura em figuras e tabelas
\definecolor{archBanco}{HTML}{1B9E77}      % banco de experiência
\definecolor{archHibrido}{HTML}{D95F02}    % híbrido
\definecolor{archGit}{HTML}{7570B3}        % git-metáfora
\definecolor{archAST}{HTML}{E7298A}        % AST-guiado
\definecolor{archGrafo}{HTML}{66A61E}      % grafo de conhecimento
\definecolor{archTexto}{HTML}{E6AB02}      % notas de texto
\definecolor{archPaginacao}{HTML}{A6761D}  % paginação OS
\definecolor{archOutro}{HTML}{666666}      % outro / survey / benchmark

% Cores semânticas para resultados
\definecolor{posGain}{HTML}{2E7D32}        % ganho positivo
\definecolor{negLoss}{HTML}{C62828}        % degradação / perda
\definecolor{neutGray}{HTML}{9E9E9E}       % neutro / sem dados
\definecolor{hlRow}{HTML}{F5F5F5}          % fundo alternado em tabelas

% =====================================================================
% Figuras e tabelas
% =====================================================================
\usepackage{graphicx}
\usepackage{booktabs}
\usepackage{tabularx}
\usepackage{longtable}
\usepackage{multirow}

% Espaçamento entre linhas de tabelas — padrão booktabs é apertado,
% 1.15 dá respiro sem desperdiçar espaço vertical
\renewcommand{\arraystretch}{1.15}

% Formatação de captions: fonte menor, label em negrito, separador
% por ponto, margem para não colidir com as bordas da tabela.
\usepackage[font=small, labelfont=bf, labelsep=period,
            margin=1cm]{caption}

% =====================================================================
% Números e unidades — formatação PT-BR automática
% =====================================================================
% siunitx v3: vírgula decimal, ponto de milhar, alinhamento em tabelas
\usepackage{siunitx}
\sisetup{
  output-decimal-marker = {,},
  group-separator = {.},
  group-minimum-digits = 4,
  detect-all,
}
% Uso: \num{1036} → 1.036 | \num{6.5} → 6,5 | \SI{86}{\percent}
% Em tabelas: coluna S para alinhamento numérico automático

% =====================================================================
% Listas compactas
% =====================================================================
\usepackage{enumitem}
\setlist{nosep, leftmargin=*}

% =====================================================================
% TikZ — bibliotecas e estilos globais
% =====================================================================
\usepackage{tikz}
\usetikzlibrary{
  arrows.meta,
  positioning,
  shapes.geometric,
  calc,
  patterns,
  fit,
  backgrounds,
}

% Estilos reutilizáveis para diagramas
\tikzset{
  % Caixas do diagrama PRISMA
  prisma box/.style={
    draw=black!70, fill=white, rounded corners=2pt,
    minimum width=4.5cm, minimum height=1cm,
    font=\small, align=center, text width=4.2cm,
  },
  prisma arrow/.style={
    -Stealth, thick, black!70,
  },
  % Estilo para nós de taxonomia / bubble chart
  taxon node/.style={
    circle, draw=#1!70, fill=#1!15, thick,
    font=\small\bfseries, inner sep=2pt,
  },
  taxon node/.default=archBanco,
  % Barras de gráfico
  bar style/.style={
    fill=#1, draw=#1!80!black, line width=0.3pt,
  },
  bar style/.default=archBanco,
  % Label de eixo
  axis label/.style={
    font=\footnotesize, text=black!80,
  },
  % Anotação
  annot/.style={
    font=\scriptsize, text=black!60,
  },
}

% =====================================================================
% Citação ABNT (venue brasileiro)
% =====================================================================
% hyperref ANTES de abntex2cite para detecção correta de ABNThyperref
\usepackage{hyperref}
\hypersetup{
  colorlinks=true,
  linkcolor=black!80!blue,
  citecolor=black!80!blue,
  urlcolor=black!80!blue,
  bookmarksdepth=3,
}
\usepackage[alf,abnt-etal-list=3]{abntex2cite}

% =====================================================================
% Layout
% =====================================================================
% Margens padrão ABNT
\usepackage[left=3cm,right=2cm,top=3cm,bottom=2cm]{geometry}

% Para facilitar adaptação futura:
% - Trocar babel para [english] para venue internacional
% - Trocar abntex2cite para natbib ou biblatex
% - Trocar template do documentclass conforme venue

\begin{document}

\title{Arquiteturas de Persistência de Conhecimento para Agentes de
Codificação Baseados em IA: Uma Revisão Sistemática da Literatura}

\author{
  Alexandre Zatti\thanks{Programa de Pós-Graduação em Computação
  Aplicada (PPGPCA), Universidade Federal da Fronteira Sul (UFFS).
  Autor correspondente.} \\
  Samuel Feitosa\thanks{PPGPCA/UFFS. Orientador.}
}

\date{}

\maketitle

\begin{abstract}
\textbf{Contexto.}
Agentes de codificação baseados em Modelos de Linguagem de Grande
Escala (LLMs) sofrem de volatilidade de conhecimento entre sessões:
a cada nova tarefa, o agente perde o contexto acumulado em
resoluções anteriores.
Diversas arquiteturas de persistência foram propostas para mitigar
essa limitação, mas o campo permanece fragmentado, sem comparação
sistemática entre alternativas nem análise cruzada de custos.
\textbf{Objetivo.}
Mapear a literatura sobre persistência de conhecimento em agentes
de codificação baseados em Inteligência Artificial (IA) e responder
a cinco questões secundárias que cobrem taxonomia, avaliação,
desempenho, custo e complexidade arquitetural.
\textbf{Método.}
Revisão sistemática conduzida conforme o PRISMA 2020.
Buscas em Scopus e arXiv, complementadas por rastreamento de
citações via Semantic Scholar, identificaram 1.194 registros
publicados entre 2023 e março de 2026.
Após remoção de duplicatas e triagem por título, resumo e texto
completo, 28 estudos foram incluídos para extração estruturada
de 37 campos.
\textbf{Resultados.}
Bancos de experiência respondem por 46\% dos estudos primários,
com o próprio LLM atuando como controlador de leitura e escrita
em 54\% dos casos e escopo temporal híbrido em 58\%.
Os 19 estudos com comparação controlada reportam ganhos universais
sobre o baseline sem persistência, com mediana de 6,8 pontos
percentuais no SWE-bench Verified e amplitude de +2,2 a +39,0
pontos percentuais.
Apesar desses ganhos, 43\% dos estudos omitem dados de custo, 32\%
documentam cenários de degradação por memória e nenhum compara três
ou mais arquiteturas no mesmo benchmark com o mesmo modelo.
\textbf{Conclusão.}
A ausência de comparações diretas e a subnotificação de custos
limitam tanto a escolha arquitetural quanto a avaliação prática.
Esses achados motivam um estudo comparativo empírico controlado,
previsto como próxima etapa desta pesquisa.
\end{abstract}

\textbf{Palavras-chave:} revisão sistemática; agentes de codificação;
persistência de conhecimento; memória para agentes; custo-eficiência.

\section{Introdução}
\label{sec:introducao}

% PRISMA Item 3: Rationale — justificativa no contexto do conhecimento existente
% PRISMA Item 4: Objectives — declaração explícita das questões abordadas

Nove dos 28 estudos sobre persistência de conhecimento em
agentes de codificação relatam cenários em que a memória
prejudica o desempenho do agente.
Esse achado torna crítico distinguir quando persistir conhecimento
é vantajoso e quando é prejudicial, decisão que precede a
comparação entre arquiteturas tecnicamente sofisticadas.
Agentes de codificação baseados em modelos de linguagem de grande
escala (LLMs) alcançaram adoção crescente em tarefas de Engenharia
de Software como resolução automatizada de issues, geração de
código e manutenção de repositórios~\cite{du2026memory}.
Sistemas como SWE-Agent, OpenHands e Agentless interagem
iterativamente com o código-fonte, executam testes e propõem
correções sem intervenção humana direta.
A capacidade desses agentes, contudo, permanece limitada por uma
restrição básica: cada sessão parte do zero.

O agente redescobre a arquitetura do projeto, repete erros já
cometidos em sessões anteriores e ignora decisões de design
previamente validadas.
Essa volatilidade de conhecimento equivale a um desenvolvedor sem
memória entre reuniões.
A persistência de conhecimento entre sessões surgiu como resposta
a essa limitação: ao registrar informações de resoluções
anteriores, o agente pode reutilizá-las nas tarefas seguintes.

Ao menos 24 sistemas de persistência foram propostos desde 2024,
todos publicados entre 2024 e o primeiro trimestre de 2026.
As arquiteturas variam de bancos de experiência que destilam
trajetórias de resolução a grafos de conhecimento, memória guiada
por AST e metáforas de controle de versão inspiradas no Git.
O campo, porém, permanece fragmentado: os estudos não convergem
para uma taxonomia unificada, e a maioria não cita trabalhos
contemporâneos sobre o mesmo problema.

Dessa fragmentação emergem lacunas que limitam tanto a pesquisa
quanto a prática.
Nenhum estudo compara três ou mais arquiteturas de persistência no
mesmo benchmark com o mesmo modelo; toda evidência de superioridade
arquitetural depende de comparações entre estudos com condições
experimentais heterogêneas.
A subnotificação de custos agrava o problema: 43\% dos estudos
omitem dados de custo, e nenhum constrói uma fronteira de Pareto
entre acurácia e custo computacional.
Tampouco existe evidência controlada sobre a relação entre
complexidade arquitetural e eficácia, de modo que não se sabe se
arquiteturas simples são suficientes ou se a sofisticação se
justifica em cenários específicos.

Trabalhos secundários existentes cobrem aspectos adjacentes sem
preencher essas lacunas; o posicionamento detalhado desta revisão
diante deles é apresentado na
\autoref{sec:trabalhos-relacionados}.

Esta revisão sistemática responde à seguinte questão principal: \textit{de que forma a literatura existente aborda a
persistência de conhecimento em agentes de codificação baseados em
IA, e que evidências empíricas existem sobre a eficácia e
custo-eficiência de diferentes arquiteturas de persistência?}
Para operacionalizar essa questão principal, cinco questões
secundárias orientam a síntese, cada uma motivada por uma lacuna
específica no estado atual do campo.

\noindent\textbf{QS1.} Quais arquiteturas de persistência foram
propostas, e como se classificam quanto a escopo temporal,
substrato representacional e política de controle?
A heterogeneidade arquitetural impede tanto a comparação direta
quanto a adoção informada por equipes que precisam escolher entre
alternativas concorrentes.
Uma taxonomia explícita revela quais escolhas de projeto dominam
o campo e quais permanecem inexploradas.

\noindent\textbf{QS2.} Quais métodos de avaliação e benchmarks são
utilizados, e quão comparáveis são os resultados?
Quando estudos diferentes adotam benchmarks, modelos base e
protocolos experimentais distintos, ganhos reportados não são
diretamente comparáveis entre si.
Esta questão verifica se há convergência metodológica e identifica
os obstáculos que impedem a síntese cross-study.

\noindent\textbf{QS3.} Quais resultados de desempenho são
reportados em relação a baselines sem persistência?
A magnitude dos ganhos determina se persistência representa uma
melhoria marginal ou substantiva sobre os agentes sem memória.
Sintetizar esses resultados expõe se o ganho se mantém em
diferentes condições experimentais ou depende de contextos
específicos.

\noindent\textbf{QS4.} Quais métricas de custo e eficiência são
reportadas, e qual a relação entre custo e eficácia?
Persistência adiciona overhead de armazenamento, recuperação e
contexto efetivo do agente.
Sem dados de custo, qualquer ganho de acurácia é incompleto para
fins de adoção prática, e a relação custo-eficácia permanece
inacessível ao leitor.

\noindent\textbf{QS5.} Que evidências existem sobre a relação
entre complexidade arquitetural e eficácia, incluindo degradação?
Sistemas mais sofisticados nem sempre superam alternativas
simples, e mais memória nem sempre produz mais eficácia.
Identificar mecanismos de degradação e padrões de complexidade
versus ganho informa decisões de projeto em cenários onde a
sofisticação não se justifica.

\begin{table}[htbp]
\centering
\caption{Mapa das questões secundárias: objetivo de cada questão e
localização da resposta no manuscrito.}
\label{tab:mapa-rqs}
\small
\begin{tabularx}{\textwidth}{l X l l}
\toprule
\textbf{QS} & \textbf{Objetivo} & \textbf{Seção} & \textbf{Tabela/Figura} \\
\midrule
QS1 & Mapear e classificar as arquiteturas de persistência &
  \autoref{sec:qs1} & \autoref{tab:arquiteturas}, \autoref{fig:taxonomia} \\
QS2 & Catalogar métodos de avaliação e benchmarks &
  \autoref{sec:qs2} & \autoref{tab:benchmarks} \\
QS3 & Sintetizar resultados de desempenho vs.\ baselines &
  \autoref{sec:qs3} & \autoref{tab:desempenho}, \autoref{fig:ganhos} \\
QS4 & Quantificar cobertura e relação custo-eficácia &
  \autoref{sec:qs4} & \autoref{tab:custo}, \autoref{fig:custo} \\
QS5 & Avaliar relação entre complexidade e eficácia &
  \autoref{sec:qs5} & \autoref{tab:degradacao} \\
\bottomrule
\end{tabularx}
\end{table}

\noindent O escopo abrange publicações de 2023 a março de 2026, com
buscas em Scopus e arXiv complementadas por rastreamento de citações.

As contribuições desta revisão são as seguintes:

\begin{enumerate}
  \item o primeiro mapeamento sistemático do campo, classificando
    24 arquiteturas de persistência para agentes de codificação
    segundo três dimensões: escopo temporal da memória, substrato
    representacional e política de leitura e escrita;
  \item a síntese quantitativa de 19 comparações controladas sobre
    eficácia e custo, com identificação de rendimentos decrescentes
    em baselines fortes e de impacto bimodal sobre o custo
    operacional;
  \item a confirmação empírica de três lacunas centrais do campo:
    ausência de comparações head-to-head entre arquiteturas no
    mesmo benchmark, subnotificação sistemática de custos por
    tarefa e cobertura insuficiente da relação entre complexidade
    arquitetural e eficácia.
\end{enumerate}

As lacunas confirmadas motivam diretamente o estudo comparativo
empírico planejado como segunda etapa desta dissertação, descrito
em mais detalhe na \autoref{sec:direcoes-futuras}.

\Needspace{14\baselineskip}
\section{Trabalhos Relacionados}
\label{sec:trabalhos-relacionados}

% PRISMA Item 3: Rationale - posicionamento contra revisões anteriores

Esta seção posiciona a presente revisão contra duas classes de
trabalhos secundários: surveys sobre memória para agentes
baseados em LLMs em geral e revisões sistemáticas sobre LLMs
aplicados a Engenharia de Software.
A primeira classe é representada principalmente pelo levantamento
recente de Du et al.; a segunda agrupa cinco SLRs publicadas
entre 2024 e 2025.
Nenhuma das duas classes adota a persistência inter-sessão como
objeto primário, nem reporta análise de custo operacional por
tarefa, que são as dimensões críticas avaliadas nesta investigação.
A \autoref{tab:slrs-previas} sintetiza a comparação por foco
temático e cobertura dessas duas dimensões.

\begin{table}[htbp]
\centering
\caption{Comparação entre trabalhos secundários relacionados e
esta revisão. As colunas \emph{Persistência?} e \emph{Custo?}
indicam, respectivamente, se o trabalho aborda explicitamente a
persistência de conhecimento entre sessões e se reporta análise
de custo operacional por tarefa.}
\label{tab:slrs-previas}
\small
\begin{tabularx}{\textwidth}{l c X c c}
\toprule
\textbf{Trabalho} & \textbf{Ano} & \textbf{Foco principal} &
  \textbf{Persistência?} & \textbf{Custo?} \\
\midrule
Du et al.~\cite{du2026memory} & 2026 &
  Memória para agentes LLM em geral &
  Sim & Não \\
Hou et al.~\cite{hou2024llmse} & 2024 &
  LLMs para Engenharia de Software (395 estudos) &
  Não & Não \\
Zhang et al.~\cite{zhang2024unifying} & 2024 &
  Modelos de linguagem para código (NLP $\cap$ ES) &
  Não & Não \\
Jin et al.~\cite{jin2024llmagents} & 2024 &
  Evolução de LLMs para agentes em ES &
  Não & Não \\
Wang et al.~\cite{wang2024llmagents} & 2024 &
% audit:ignore-next-line — describes cited work's own framing
  Landscape de agentes autônomos para ES &
  Parcial & Não \\
Liu et al.~\cite{liu2024llmagents} & 2024 &
  Agentes LLM para Engenharia de Software (124 estudos) &
  Parcial & Não \\
\midrule
\textbf{Esta revisão} & 2026 &
  \textbf{Persistência de conhecimento em agentes de codificação
  baseados em IA (28 estudos)} &
  \textbf{Sim} & \textbf{Sim} \\
\bottomrule
\end{tabularx}
\end{table}

\subsection{Surveys sobre memória para agentes LLM}
\label{sec:rel-memoria}

\citeonline{du2026memory} apresentam o levantamento mais recente
sobre mecanismos de memória para agentes baseados em LLMs,
propondo uma taxonomia tridimensional (escopo temporal, substrato
representacional e política de controle) que esta revisão adota e
refina como ponto de partida da QS1.
O survey cataloga cinco famílias de mecanismos de memória:
compressão de contexto residente, stores aumentados por
recuperação, memória reflexiva auto-corretiva, contextos
hierárquicos virtuais e gestão de memória aprendida por política.
O escopo dos autores é explicitamente amplo, incluindo agentes
para tarefas gerais como navegação web, raciocínio matemático e
diálogo, sem restrição ao domínio de codificação.
Como consequência, a cobertura específica para Engenharia de
Software no survey de Du et al.\ limita-se a ChatDev e MetaGPT,
enquanto sistemas centrais para a presente investigação, como
MemCoder~\cite{deng2026memcoder},
CodeMEM~\cite{wang2026codemem} e a abordagem de Repository
Memory de~\citeonline{wang2026improving}, não são mencionados.
A presente revisão complementa esse trabalho restringindo o
escopo a agentes de codificação e incorporando análise de custo
operacional por tarefa, dimensão ausente no survey de Du et al.

\subsection{Revisões sistemáticas sobre LLMs em Engenharia de
Software}
\label{sec:rel-llmse}

Cinco revisões sistemáticas recentes mapeiam aplicações de LLMs em
Engenharia de Software com escopo amplo.
\citeonline{hou2024llmse} cobrem 395 estudos publicados até 2023,
organizados por tarefa de engenharia (geração de código, reparo,
teste, manutenção) e por tipo de modelo, sem subseção dedicada a
memória ou persistência cross-session.
\citeonline{zhang2024unifying} unificam perspectivas de
Processamento de Linguagem Natural e Engenharia de Software sobre
modelos para código, com ênfase em pré-treinamento e benchmarks de
avaliação.
\citeonline{jin2024llmagents} traçam a transição de LLMs para
agentes baseados em LLMs em Engenharia de Software, mas tratam
memória como componente arquitetural genérico, sem caracterizar
arquiteturas de persistência inter-sessão.
\citeonline{wang2024llmagents} mapeiam o cenário dos agentes
autônomos para Engenharia de Software com foco em pipelines de
execução, e \citeonline{liu2024llmagents} cobrem 124 estudos com
discussão parcial sobre módulos de memória, porém sem síntese
cross-system de eficácia e custo da persistência.
Nenhuma dessas revisões responde diretamente às questões
secundárias QS1--QS5 desta investigação: a taxonomia de
arquiteturas de persistência, a comparabilidade de benchmarks
específicos para memória, a quantificação de ganhos sobre
baselines sem memória, a relação custo-eficácia ou os mecanismos
de degradação por memória.

\subsection{Lacuna posicional desta revisão}
\label{sec:rel-lacuna}

Esta revisão restringe o escopo aos agentes de codificação
baseados em IA e adota a persistência de conhecimento entre
sessões como objeto primário de investigação, não como componente
arquitetural acessório de uma síntese mais ampla.
Sobre essa base, incorpora a análise de custo operacional por
tarefa como dimensão dedicada (QS4), o que constitui o
diferencial central frente aos trabalhos anteriores: nenhuma das
SLRs prévias identificadas reúne simultaneamente foco em
persistência inter-sessão, escopo de agentes de codificação e
quantificação cross-study de custos.
Essa configuração permite responder questões que os trabalhos
secundários anteriores não abordam, sobretudo a ausência de
comparações head-to-head entre arquiteturas, a subnotificação
sistemática de custos e os mecanismos pelos quais a memória pode
prejudicar o desempenho do agente.

\Needspace{14\baselineskip}
\section{Método}
\label{sec:metodo}

Esta revisão sistemática segue a metodologia de
\citeonline{kitchenham2007guidelines} para revisões sistemáticas em
Engenharia de Software, incorporando as atualizações de
\citeonline{kitchenham2013systematic}.
O reporte adota as Software Engineering Guidelines for Reporting
Secondary Studies (SEGRESS)~\cite{kitchenham2023segress}, adaptação do
PRISMA 2020~\cite{page2021prisma} para estudos secundários em Engenharia
de Software.
A estratégia de busca segue a extensão
PRISMA-S~\cite{rethlefsen2021prismas}, e o reporte completo segue o
checklist de 27 itens do PRISMA 2020.

\subsection{Protocolo e registro}
\label{sec:protocolo}

O protocolo foi elaborado antes do início das buscas.
Seis emendas foram registradas durante a execução, refinando critérios
de elegibilidade e o formulário de extração.

\subsection{Critérios de elegibilidade}
\label{sec:criterios}

A formulação dos critérios seguiu o \textit{framework} PICo:
a \textbf{P}opulação compreende agentes de codificação baseados em IA
com mecanismos de persistência de conhecimento entre sessões;
o \textbf{I}nteresse abrange as arquiteturas de persistência, sua
eficácia e custo-eficiência;
o \textbf{Co}ntexto inclui avaliações empíricas, benchmarks e propostas
de ferramentas ou frameworks.

A \autoref{tab:criterios} apresenta os critérios de inclusão e exclusão.
CI1 exige componente explícito de armazenamento e reutilização de
conhecimento entre sessões ou tarefas sequenciais;
sistemas que apenas recuperam contexto dentro de uma única sessão
não atendem ao critério.
CI2 estabelece piso mínimo de avaliação empírica com pelo menos duas
tarefas e métricas quantitativas.
CE2 opera como regra condicional: papers sobre agentes LLM genéricos
foram incluídos quando a avaliação continha benchmarks de Engenharia
de Software.
O período de cobertura abrange publicações de 2023 a março de 2026.

\begin{table}[htbp]
  \centering
  \caption{Critérios de elegibilidade.}
  \label{tab:criterios}
  \begin{tabularx}{\textwidth}{l l X}
    \toprule
    ID  & Tipo & Descrição \\
    \midrule
    CI1 & Inclusão & Propõe ou avalia mecanismo para persistir
      conhecimento entre sessões em agentes de codificação baseados
      em IA \\
    CI2 & Inclusão & Reporta resultados empíricos com pelo menos
      duas tarefas e métricas quantitativas \\
    CI3 & Inclusão & Publicado entre 2023 e março de 2026 \\
    CI4 & Inclusão & Escrito em inglês ou português \\
    CI5 & Inclusão & Publicação com revisão por pares ou preprint com
      autores e instituições identificáveis \\
    \midrule
    CE1 & Exclusão & Aborda apenas RAG ou recuperação de contexto
      intra-sessão, sem persistência entre sessões \\
    CE2 & Exclusão & Aborda apenas tarefas de linguagem natural, sem
      avaliação em benchmarks de Engenharia de Software \\
    CE3 & Exclusão & Avaliação limitada a exemplos anedóticos \\
    CE4 & Exclusão & Tutorial, editorial ou position paper sem
      contribuição técnica \\
    CE5 & Exclusão & Versão duplicada ou supersedida por versão posterior
      dos mesmos autores \\
    \bottomrule
  \end{tabularx}
\end{table}

\subsection{Fontes de informação}
\label{sec:fontes}

As buscas foram realizadas em duas bases de dados com suporte a busca
booleana reproduzível: Scopus (via API) e arXiv (via API REST com
paginação completa).
A estratégia de duas bases complementada por rastreamento de citações é
validada empiricamente por \citeonline{mourao2020hybrid}, que demonstraram
que Scopus combinado com \textit{snowballing} alcança \textit{recall} comparável a buscas em
quatro ou mais bases para revisões sistemáticas em Engenharia de Software.
O IEEE Xplore não foi buscado separadamente porque o Scopus indexa
integralmente seus proceedings e journals.

A busca foi suplementada por rastreamento de citações forward e backward
via Semantic Scholar a partir de sete artigos-semente selecionados por
representarem diversidade arquitetural e serem os mais citados no
subcampo~\cite{du2026memory,joshi2025swebenchcl,wang2026codemem,deng2026memcoder,wu2025gcc,wang2026improving,lindenbauer2025ctimrover}.
Um dos sete artigos-semente (CodeMEM) não foi capturado pelas strings de
busca por utilizar terminologia fora do vocabulário das strings;
sua adição manual ao \textit{pool} de triagem foi registrada como emenda ao
protocolo.

\subsection{Estratégia de busca}
\label{sec:busca}

A string de busca combina três blocos conceituais conectados por AND:
um bloco sobre agentes de IA para código (C1, 12 termos) com um segundo
sobre memória e persistência (C2, 13 termos), restritos ao domínio de
engenharia de software por um terceiro bloco (C3, 8 termos).
A string foi adaptada para a sintaxe de cada base: TITLE-ABS-KEY no
Scopus com filtro \texttt{PUBYEAR\,>\,2022}, e campos \texttt{abs:}/\texttt{ti:}
no arXiv com filtro de categorias cs.SE, cs.AI e cs.CL.\footnote{String
Scopus: \texttt{TITLE-ABS-KEY(("LLM agent" OR "AI coding agent" OR
"software engineering agent" OR "autonomous coding" OR "code agent" OR
"LLM-based software" OR "AI-assisted software" OR "coding assistant" OR
"language agent" OR "code assistant" OR "agentic coding" OR "SWE agent")
AND ("memory" OR "knowledge persistence" OR "experience" OR "context
management" OR "long-term memory" OR "episodic memory" OR "knowledge
reuse" OR "continual learning" OR "experience replay" OR
"memory-augmented" OR "persistent memory" OR "persistent context" OR
"cross-task") AND ("software engineering" OR "software development" OR
"code generation" OR "bug fix" OR "code repair" OR "repository" OR
"software maintenance" OR "issue resolution")) AND PUBYEAR > 2022}.
A string completa para o arXiv e o checklist PRISMA-S estão disponíveis
no material suplementar.}

Antes da execução, as strings passaram por auditoria contra oito SLRs
comparáveis, 16 sistemas conhecidos e dez formulações alternativas de
termos, resultando na adição de seis termos e cobertura de seis dos sete
artigos-semente.
O único artigo-semente não recuperado (CodeMEM) utiliza terminologia fora
do vocabulário das strings e foi adicionado manualmente.

A \autoref{tab:busca} apresenta os resultados por fonte.
As buscas foram executadas em 25 de março de 2026.
O rastreamento de citações identificou 186 registros únicos (8 forward,
178 backward).
Após consolidação e desduplicação, o \textit{pool} inicial contém 1.111 registros
para triagem.

\begin{table}[htbp]
  \centering
  \caption{Resultados das buscas por fonte.}
  \label{tab:busca}
  \begin{tabular}{l r}
    \toprule
    Fonte & Registros \\
    \midrule
    Scopus            &    39 \\
    arXiv             &   968 \\
    Citações forward  &     8 \\
    Citações backward &   178 \\
    Manual            &     1 \\
    \midrule
    Total bruto       & 1.194 \\
    Duplicatas removidas & $-83$ \\
    \textbf{Total para triagem} & \textbf{1.111} \\
    \bottomrule
  \end{tabular}
\end{table}

\subsection{Processo de seleção}
\label{sec:selecao}

A seleção ocorreu em duas fases.
Os 1.111 registros foram triados por título e resumo pelo
pesquisador, aplicando os critérios CI/CE.
Registros classificados como incertos avançaram para a fase seguinte.
Dos 1.111 triados, 1.060 foram excluídos e 51 avançaram para avaliação
por texto completo.

A leitura integral dos 51 artigos resultou em 23 exclusões.
Dos 23 excluídos, 12 foram removidos por CE2 (agentes LLM sem avaliação
em Engenharia de Software) e 11 pela combinação CI1/CE1 (recuperação
intra-sessão sem persistência entre sessões).
Três casos limítrofes foram aceitos com justificativa documentada.
Ao final, 28 estudos foram incluídos para extração de dados.

\subsection{Extração de dados}
\label{sec:extracao}

Os dados foram extraídos com formulário de 37 campos organizados em
sete categorias, com ênfase em arquitetura de persistência (8 campos)
e custo e eficiência (6 campos), que correspondem diretamente às
questões de pesquisa.
Dezesseis campos utilizam vocabulário controlado com valores predefinidos.
Campos sem dados disponíveis recebem NR (não reportado) quando o paper não
menciona o dado ou NA (não aplicável) quando o campo não se aplica ao
sistema, conforme convenção registrada como emenda ao protocolo.

Cada paper foi lido, seus dados registrados no formulário e os valores
verificados em seguida.
A validação incluiu verificação programática de 448 valores enum e
\textit{spot-check} de 40 campos interpretativos em dez papers.
A taxa de erro final foi inferior a 1\%, correspondente a oito
correções em 1.036 valores extraídos (28 papers $\times$ 37 campos).

\subsection{Avaliação de qualidade}
\label{sec:qualidade}

A qualidade dos estudos foi avaliada com checklist de seis itens
adaptado de \citeonline{kitchenham2007guidelines}: (Q1) objetivos
claramente definidos, (Q2) método de avaliação adequado, (Q3) ameaças
à validade discutidas, (Q4) resultados baseados em dados empíricos,
(Q5) contexto descrito para replicação e (Q6) métricas claramente
definidas.
Cada item recebeu pontuação 0, 0,5 ou 1, com máximo de 6,0.

A distribuição das pontuações varia de 2,0 a 6,0, com mediana de 5,0.
Vinte e seis dos 28 estudos (93\%) obtiveram pontuação igual ou
superior a 4,0.
Os dois estudos com pontuação inferior a
3,0~\cite{bui2026opendev,du2026memory} contribuem apenas descrições
qualitativas e não impactam os achados quantitativos, conforme
verificado na análise de sensibilidade (\autoref{sec:resultados}).
Para calibração, oito papers foram reavaliados após a extração completa.

\subsection{Métodos de síntese}
\label{sec:sintese}

A síntese é narrativa, organizada tematicamente pelas cinco questões
secundárias (QS1 a QS5).
Meta-análise foi descartada pela heterogeneidade de benchmarks, modelos
base e condições experimentais entre os estudos.

Duas análises de sensibilidade foram conduzidas:
(S1) re-síntese restrita aos cinco estudos publicados em venues com
revisão por pares, para verificar se os achados se sustentam sem
preprints;
(S2) exclusão dos dois estudos com pontuação inferior a 3,0, avaliando
o impacto nos achados quantitativos.

A avaliação de viés de publicação formal não é aplicável para síntese
narrativa~\cite{page2021prisma}.
O predomínio de preprints no corpus (82\%) é reconhecido como fonte
potencial de viés e discutido nas limitações.
A certeza da evidência foi avaliada qualitativamente por achado:
alta (múltiplos estudos de qualidade $\geq$~4/6 convergem),
moderada (poucos estudos ou qualidade variável) e
baixa (achado baseado em estudo único ou de baixa
qualidade)~\cite{kitchenham2023segress}.

\Needspace{14\baselineskip}
\section{Resultados}
\label{sec:resultados}

\subsection{Seleção de estudos}
\label{sec:selecao-estudos}

A \autoref{fig:prisma-flow} apresenta o diagrama de fluxo PRISMA.
As buscas em Scopus e arXiv, complementadas por rastreamento de citações
e uma adição manual, identificaram 1.194 registros.
Após remoção de 83 duplicatas, 1.111 registros foram triados por título
e resumo.
Dessa triagem, 1.060 foram excluídos e 51 avançaram para leitura de
texto completo.

Das 23 exclusões na fase de texto completo, 12 ocorreram por CE2
(agentes LLM sem avaliação em Engenharia de Software) e 11 pela
combinação CI1/CE1 (recuperação intra-sessão sem persistência entre
sessões).
Três casos limítrofes foram aceitos com justificativa documentada.
A lista completa dos 23 estudos excluídos na triagem de texto completo,
com as respectivas razões de exclusão por critério, está disponível no
pacote de replicação (DOI: 10.5281/zenodo.19463131).
Ao final, 28 estudos foram incluídos para extração de dados.

% O diagrama PRISMA já contém \begin{figure} e \label{fig:prisma-flow}
% Requer: tikz, calc
\begin{figure}[!htbp]
\centering
\resizebox{\textwidth}{!}{%
\begin{tikzpicture}[
  box/.style={rectangle, draw=black, thick, text width=5.5cm,
    minimum height=1.2cm, align=center, font=\small},
  boxwide/.style={rectangle, draw=black, thick, text width=11.0cm,
    minimum height=1.0cm, align=center, font=\small},
  boxgray/.style={rectangle, draw=black, thick, fill=gray!15,
    text width=5.5cm, minimum height=1.2cm, align=center, font=\small},
  phaselabel/.style={rectangle, draw=black, thick, fill=blue!10,
    minimum height=1.0cm, align=center,
    font=\small\bfseries, rotate=90, inner sep=3pt},
  arrow/.style={->, thick, >=stealth},
]
\node[phaselabel, minimum width=3.0cm] (l1) at (-2, -0.5)
  {Identificação};
\node[phaselabel, minimum width=6.0cm] (l2) at (-2, -7.0)
  {Triagem};
\node[phaselabel, minimum width=3.5cm] (l3) at (-2, -13.0)
  {Inclusão};

\node[box] (db) at (2, 0)
  {Registros identificados em\\bases de dados ($n = 1.007$):\\
   Scopus ($n = 39$)\\arXiv ($n = 968$)};

\node[box] (other) at (11, 0)
  {Registros identificados por\\outros métodos ($n = 187$):\\
   Rastreamento de citações ($n = 186$)\\
   Adição manual ($n = 1$)};

\node[boxwide] (dedup) at (6.5, -2.5)
  {Registros após remoção de duplicatas ($n = 1.111$)\\
   {\footnotesize Total identificado: $n = 1.194$\quad
    Duplicatas removidas: $n = 83$}};

\draw[arrow] (db) -- (dedup);
\draw[arrow] (other) -- (dedup);

\node[box] (ta) at (2, -5.0)
  {Registros triados por\\título e resumo\\($n = 1.111$)};

\node[boxgray] (ta_excl) at (11, -5.0)
  {Registros excluídos\\($n = 1.060$)};

\draw[arrow] (dedup) -- (ta);
\draw[arrow] (ta) -- (ta_excl);

\node[boxgray] (nr) at (11, -7.0)
  {Relatórios não recuperados\\($n = 0$)};

\node[box] (ft) at (2, -9.0)
  {Relatórios avaliados por\\texto completo\\($n = 51$)};

\node[boxgray, minimum height=2.0cm] (ft_excl) at (11, -9.2)
  {Relatórios excluídos ($n = 23$):\\[3pt]
   {\footnotesize
   CE2 (sem avaliação em Eng.\ Soft.): $n = 12$\\[2pt]
   CI1/CE1 (sem persistência): $n = 11$}};

\draw[arrow] (ta) -- (ft);
\draw[arrow] (2, -7.0) -- (nr.west);
\draw[arrow] (ft) -- (ft_excl);

\node[box, fill=green!10] (incl) at (2, -11.8)
  {Estudos incluídos na\\revisão ($n = 28$)};

\draw[arrow] (ft) -- (incl);

\node[box, fill=green!10] (qual) at (2, -14.0)
  {Estudos com dados extraídos\\e avaliação de qualidade\\($n = 28$)};

\draw[arrow] (incl) -- (qual);

\end{tikzpicture}%
}
\caption{Diagrama de fluxo PRISMA 2020 da revisão sistemática.
Adaptado de \citeonline{page2021prisma}.}
\label{fig:prisma-flow}
\end{figure}


\subsection{Características dos estudos incluídos}
\label{sec:caracteristicas}

A \autoref{tab:caracteristicas} sumariza os 28 estudos incluídos.
Todos foram publicados entre 2024 e 2026; nenhum estudo de 2023
atendeu aos critérios de elegibilidade.
Preprints representam 82\% do corpus (23/28), com duas publicações em
conferência~\cite{wang2026improving,zhang2025darwingodel} e três em
workshop~\cite{deshpande2025memtrack,shen2025talm,lindenbauer2025ctimrover}.

Quanto ao código-fonte, 17 dos 24 estudos primários disponibilizam repositório aberto e
três oferecem acesso parcial.
Os frameworks de agente mais recorrentes são implementações customizadas,
seguidas por SWE-Agent e OpenHands.
Os LLMs utilizados incluem GPT-4o, Claude 3.5/4/4.5 Sonnet,
DeepSeek-V3 e Gemini 2.5/3 Pro.

Quatro estudos não propõem arquiteturas primárias de persistência:
\citeonline{du2026memory} é survey,
\citeonline{deshpande2025memtrack} e \citeonline{joshi2025swebenchcl}
propõem benchmarks, e
\citeonline{srivastava2025memorygraft} investiga segurança de memória.
Os 24 estudos restantes propõem ou implementam arquiteturas e são
analisados nas subseções seguintes.

\begin{longtable}{l c l l l S[table-format=1.1]}
  \caption{Características dos 28 estudos incluídos.
    Score: pontuação de qualidade (escala 0--6, conforme
    \autoref{sec:qualidade-resultados}).}
  \label{tab:caracteristicas} \\
  \toprule
  Estudo & Ano & Venue & Tipo & Arquitetura & {Score} \\
  \midrule
  \endfirsthead
  \toprule
  Estudo & Ano & Venue & Tipo & Arquitetura & {Score} \\
  \midrule
  \endhead
  \bottomrule
  \endfoot
  Bui (OpenDev) & 2026 & arXiv & preprint & notas texto & 2.0 \\
  Chen (SWE-Exp) & 2025 & arXiv & preprint & banco exp. & 6.0 \\
  Deng (MemCoder) & 2026 & arXiv & preprint & banco exp. & 4.5 \\
  Deshpande (MemTrack) & 2025 & NeurIPS-W & workshop & benchmark & 4.5 \\
  Dong (KernelBlaster) & 2026 & arXiv & preprint & grafo conh. & 4.5 \\
  Du (survey) & 2026 & arXiv & preprint & survey & 2.5 \\
  Guo (EET) & 2026 & arXiv & preprint & banco exp. & 5.5 \\
  Hosain (Xolver) & 2025 & arXiv & preprint & híbrido & 5.5 \\
  Joshi (SWE-Bench-CL) & 2025 & arXiv & preprint & benchmark & 5.0 \\
  Li (GBT) & 2026 & arXiv & preprint & outro & 5.5 \\
  Lindenbauer (CTIM-Rover) & 2025 & REALM-W & workshop & banco exp. & 4.0 \\
  Nadafian (KAPSO) & 2026 & arXiv & preprint & híbrido & 5.5 \\
  Qian (IER) & 2024 & arXiv & preprint & banco exp. & 4.0 \\
  Shen (TALM) & 2025 & AGENT-W & workshop & banco exp. & 4.5 \\
  Shen (Structurally) & 2026 & arXiv & preprint & banco exp. & 5.0 \\
  Srivastava (MemoryGraft) & 2025 & arXiv & preprint & segurança & 4.0 \\
  Su (Learn-by-Interact) & 2025 & arXiv & preprint & banco exp. & 5.0 \\
  Tablan (Spark) & 2025 & arXiv & preprint & banco exp. & 5.0 \\
  Vasilopoulos (Codified) & 2026 & arXiv & preprint & híbrido & 5.0 \\
  Wang (RepoMem) & 2026 & ICLR & conferência & git metáf. & 5.0 \\
  Wang (CodeMEM) & 2026 & arXiv & preprint & AST guiado & 5.0 \\
  Wang (MemGovern) & 2026 & arXiv & preprint & banco exp. & 5.5 \\
  Wong (CCA) & 2025 & arXiv & preprint & híbrido & 5.5 \\
  Wu (GCC) & 2025 & arXiv & preprint & git metáf. & 5.5 \\
  Xia (Live-SWE-Agent) & 2025 & arXiv & preprint & outro & 5.5 \\
  Zhang (AccelOpt) & 2025 & arXiv & preprint & banco exp. & 5.5 \\
  Zhang (Darwin-Gödel) & 2025 & ICLR & conferência & outro & 5.0 \\
  Zhou (ToM-SWE) & 2025 & arXiv & preprint & híbrido & 5.5 \\
\end{longtable}

A \autoref{fig:mapa-campo} oferece uma visão cronológica complementar
do corpus, distribuindo os 28 estudos por ano de publicação e
família arquitetural, com destaque para os sete artigos-semente e
para os três sistemas cujo substrato de persistência é executável.

% Figura: mapa do campo — 28 estudos por ano e arquitetura
% Eixo X: ano de publicação (2024–2026).
% Cor: família arquitetural (paleta global archXxx do main.tex).
% Marcadores: borda dupla = artigo-semente; quadrado = substrato
% de persistência executável.
% Coordenadas explícitas (sem \foreach, sem \begin{scope}[shift=...])
% para passar no verificador estático de colisões TikZ.
\begin{figure}[!htbp]
  \centering
  \begin{tikzpicture}[
    every node/.style={font=\scriptsize},
    fnode/.style={
      circle, draw=#1!75!black, fill=#1!18, line width=0.4pt,
      minimum size=10mm, inner sep=0pt,
      font=\scriptsize, text=black, align=center,
    },
    fnode/.default=archBanco,
    seednode/.style={
      circle, draw=#1!80!black, fill=#1!22,
      double, double distance=0.5pt, line width=0.5pt,
      minimum size=10mm, inner sep=0pt,
      font=\scriptsize\bfseries, text=black, align=center,
    },
    seednode/.default=archBanco,
    execnode/.style={
      rectangle, rounded corners=1pt,
      draw=#1!80!black, fill=#1!22, line width=0.5pt,
      minimum width=10mm, minimum height=10mm, inner sep=0pt,
      font=\scriptsize, text=black, align=center,
    },
    execnode/.default=archOutro,
    % (swatches usam minimum size=4mm inline para que o verificador
    % est\'atico de colis\~oes leia as dimens\~oes reais)
  ]

    % --- R\'otulos de ano no topo ---
    \node[font=\footnotesize\bfseries] at (1.50, 10.90) {2024};
    \node[font=\footnotesize\bfseries] at (5.75, 10.90) {2025};
    \node[font=\footnotesize\bfseries] at (11.25, 10.90) {2026};

    % =================================================================
    % 2024 (1 paper)
    % Coluna em x=1.50; linha em y=10.00.
    % =================================================================
    \node[fnode=archBanco] (n01) at (1.50, 10.00) {Qian};

    % =================================================================
    % 2025 (16 papers) — 3 sub-colunas em x=4.50, 5.75, 7.00
    % Linhas em y = 10.00, 8.80, 7.60, 6.40, 5.20, 4.00, 2.80
    %   (passo 1.20 cm, gap entre n\'os 0.20 cm = MAJOR-safe).
    % =================================================================
    % --- sub-coluna A ---
    \node[fnode=archBanco]   (n02) at (4.50, 10.00) {Chen};
    \node[fnode=archOutro]   (n03) at (4.50,  8.80) {Desh.};
    \node[fnode=archHibrido] (n04) at (4.50,  7.60) {Hosain};
    \node[seednode=archOutro](n05) at (4.50,  6.40) {Joshi};
    \node[seednode=archBanco](n06) at (4.50,  5.20) {Lind.};
    \node[fnode=archOutro]   (n07) at (4.50,  4.00) {Sriva.};
    \node[fnode=archBanco]   (n08) at (4.50,  2.80) {Su};

    % --- sub-coluna B ---
    \node[fnode=archBanco]   (n09) at (5.75, 10.00) {Shen};
    \node[fnode=archBanco]   (n10) at (5.75,  8.80) {Tab.};
    \node[seednode=archGit]  (n11) at (5.75,  7.60) {Repo};
    \node[fnode=archHibrido] (n12) at (5.75,  6.40) {Wong};
    \node[seednode=archGit]  (n13) at (5.75,  5.20) {Wu};
    \node[execnode=archOutro](n14) at (5.75,  4.00) {Live};
    \node[fnode=archBanco]   (n15) at (5.75,  2.80) {Accel};

    % --- sub-coluna C ---
    \node[execnode=archOutro](n16) at (7.00, 10.00) {Darw.};
    \node[fnode=archHibrido] (n17) at (7.00,  8.80) {ToM};

    % =================================================================
    % 2026 (11 papers) — 2 sub-colunas em x=10.00 e x=11.25, e 1 em 12.50
    % =================================================================
    % --- sub-coluna A ---
    \node[fnode=archTexto]   (n18) at (10.00, 10.00) {Bui};
    \node[seednode=archBanco](n19) at (10.00,  8.80) {Deng};
    \node[fnode=archGrafo]   (n20) at (10.00,  7.60) {Dong};
    \node[seednode=archOutro](n21) at (10.00,  6.40) {Du};
    \node[fnode=archBanco]   (n22) at (10.00,  5.20) {Guo};
    \node[execnode=archOutro](n23) at (10.00,  4.00) {Li};

    % --- sub-coluna B ---
    \node[fnode=archHibrido] (n24) at (11.25, 10.00) {Nadaf.};
    \node[fnode=archBanco]   (n25) at (11.25,  8.80) {Shen};
    \node[fnode=archHibrido] (n26) at (11.25,  7.60) {Vasil.};
    \node[seednode=archAST]  (n27) at (11.25,  6.40) {Wang};
    \node[fnode=archBanco]   (n28) at (11.25,  5.20) {MemG};

    % =================================================================
    % R\'otulo de eixo (abaixo da \'area do gr\'afico)
    % =================================================================
    \node[font=\footnotesize, text=black!70]
      at (6.00, 1.60) {Ano de publica\c{c}\~ao};

    % =================================================================
    % Legenda — duas linhas, coordenadas absolutas (sem scope shift)
    % Linha superior: 4 swatches de cor (y = 0.40)
    % Linha inferior: 4 swatches de cor + marcadores (y = -0.40)
    % Espa\c{c}amento horizontal entre slots = 3.40 cm
    % =================================================================
    % --- Cabe\c{c}alho da legenda ---
    \node[font=\scriptsize\bfseries, anchor=west]
      at (0.00, 1.05) {Fam\'ilia arquitetural};

    % --- Linha 1 (slots em x = 0.20, 3.80, 7.40, 11.00) ---
    \node[circle, draw=archBanco!75!black, fill=archBanco!18,
          minimum size=4mm, inner sep=0pt]
      at (0.20, 0.40) {};
    \node[anchor=west, font=\scriptsize] at (0.80, 0.40)
      {banco exp.};
    \node[circle, draw=archHibrido!75!black, fill=archHibrido!18,
          minimum size=4mm, inner sep=0pt]
      at (3.80, 0.40) {};
    \node[anchor=west, font=\scriptsize] at (4.40, 0.40)
      {h\'ibrido};
    \node[circle, draw=archGit!75!black, fill=archGit!18,
          minimum size=4mm, inner sep=0pt]
      at (7.40, 0.40) {};
    \node[anchor=west, font=\scriptsize] at (8.00, 0.40)
      {git-met\'afora};
    \node[circle, draw=archAST!75!black, fill=archAST!18,
          minimum size=4mm, inner sep=0pt]
      at (11.00, 0.40) {};
    \node[anchor=west, font=\scriptsize] at (11.60, 0.40)
      {AST-guiado};

    % --- Linha 2 ---
    \node[circle, draw=archGrafo!75!black, fill=archGrafo!18,
          minimum size=4mm, inner sep=0pt]
      at (0.20, -0.40) {};
    \node[anchor=west, font=\scriptsize] at (0.80, -0.40)
      {grafo conhec.};
    \node[circle, draw=archTexto!75!black, fill=archTexto!18,
          minimum size=4mm, inner sep=0pt]
      at (3.80, -0.40) {};
    \node[anchor=west, font=\scriptsize] at (4.40, -0.40)
      {notas texto};
    \node[circle, draw=archPaginacao!75!black, fill=archPaginacao!18,
          minimum size=4mm, inner sep=0pt]
      at (7.40, -0.40) {};
    \node[anchor=west, font=\scriptsize] at (8.00, -0.40)
      {pagina\c{c}\~ao OS};
    \node[circle, draw=archOutro!75!black, fill=archOutro!18,
          minimum size=4mm, inner sep=0pt]
      at (11.00, -0.40) {};
    \node[anchor=west, font=\scriptsize] at (11.60, -0.40)
      {outro/survey};

    % --- Linha 3: marcadores ---
    \node[font=\scriptsize\bfseries, anchor=west]
      at (0.00, -1.20) {Marcadores};
    \node[circle, draw=black!70, fill=black!8, double,
          double distance=0.5pt, line width=0.4pt,
          minimum size=4mm, inner sep=0pt]
      at (2.40, -1.20) {};
    \node[anchor=west, font=\scriptsize] at (3.00, -1.20)
      {artigo-semente (n=7)};
    \node[rectangle, rounded corners=0.5pt, draw=black!70,
          fill=black!8, line width=0.4pt,
          minimum width=4mm, minimum height=4mm, inner sep=0pt]
      at (7.40, -1.20) {};
    \node[anchor=west, font=\scriptsize] at (8.00, -1.20)
      {substrato execut\'avel (n=3)};

  \end{tikzpicture}
  \caption[Mapa cronol\'ogico do campo]{Mapa cronol\'ogico das 28
  arquiteturas inclu\'idas, distribu\'idas por ano de publica\c{c}\~ao
  (2024--2026) e fam\'ilia arquitetural (cor). Marcadores com borda
  dupla e r\'otulo em negrito identificam os sete artigos-semente;
  marcadores quadrados, os tr\^es estudos cujo substrato de
  persist\^encia \'e execut\'avel (Live SWE-Agent, Darwin-G\"odel
  Machine, Traversal-as-Policy). Fam\'ilias: bancos de
  experi\^encia (verde-azulado), h\'ibridos (laranja), git-met\'afora
  (roxo), AST-guiado (rosa), grafo de conhecimento (verde), notas
  de texto (amarelo), pagina\c{c}\~ao OS (marrom), outros (cinza).}
  \label{fig:mapa-campo}
\end{figure}


\subsection{Avaliação de qualidade}
\label{sec:qualidade-resultados}

As pontuações de qualidade variam de 2,0 a 6,0, com mediana de 5,0.
Vinte e seis dos 28 estudos (93\%) obtiveram pontuação igual ou
superior a 4,0.
Os dois estudos com pontuação inferior a
3,0~\cite{bui2026opendev,du2026memory} contribuem apenas descrições
qualitativas; sua exclusão na análise de sensibilidade S2 não altera
nenhum achado quantitativo.

\subsection{Taxonomia de arquiteturas de persistência (QS1)}
\label{sec:qs1}

A \autoref{tab:arquiteturas} apresenta a distribuição dos 24 estudos
primários por tipo de arquitetura.
Onze estudos (46\%) destilam trajetórias de resolução em entradas
reutilizáveis, tornando bancos de experiência o tipo mais frequente
(e.g.,~\cite{chen2025sweexp,deng2026memcoder,guo2026eet,shen2026structurally,su2025learnbyinteract}).
Arquiteturas híbridas respondem por cinco estudos (21\%), tipicamente
combinando memória episódica com
semântica~\cite{hosain2025xolver,nadafian2026kapso,wong2025confucius}.
Duas implementações adotam metáfora
git~\cite{wang2026improving,wu2025gcc}, e categorias como AST
guiado~\cite{wang2026codemem} e grafo de
conhecimento~\cite{dong2026kernelblaster} contam com um estudo cada.

\begin{table}[htbp]
  \centering
  \caption{Distribuição dos estudos primários por tipo de arquitetura.}
  \label{tab:arquiteturas}
  \begin{tabular}{l r r}
    \toprule
    Tipo de arquitetura & $N$ & \% \\
    \midrule
    Banco de experiência   & 11 & 46 \\
    Híbrido                &  5 & 21 \\
    Outro                  &  3 & 12 \\
    Git metáfora           &  2 &  8 \\
    AST guiado             &  1 &  4 \\
    Grafo de conhecimento  &  1 &  4 \\
    Notas de texto         &  1 &  4 \\
    \midrule
    \textbf{Total}         & \textbf{24} & \\
    \bottomrule
  \end{tabular}
\end{table}

A classificação tridimensional proposta por
\citeonline{du2026memory} permite análise mais granular.
Quanto ao escopo temporal, 14 estudos (58\%) adotam escopo híbrido,
combinando memória de curto e longo prazo; cinco (21\%) adotam escopo
procedural e quatro (17\%) escopo episódico.
Apenas \citeonline{wang2026memgovern} utiliza escopo puramente semântico.
O substrato representacional mais frequente é texto em contexto (9
estudos, 38\%), seguido por estruturado (7, 29\%) e híbrido (5, 21\%).
Três estudos~\cite{li2026traversalaspolicy,xia2025livesweagent,zhang2025darwingodel} utilizam substratos executáveis, onde a memória
assume a forma de árvores de comportamento, scripts ou agentes inteiros,
contrastando com o padrão predominante de dados passivos para
recuperação.
Quanto à política de controle, 13 estudos (54\%) delegam ao próprio LLM
as decisões de escrita e leitura da memória, nove (38\%) empregam
heurísticas fixas e dois (8\%) adotam políticas aprendidas.

A recuperação híbrida, combinando BM25, embeddings e heurísticas, é o
método dominante (11 estudos, 46\%).
Carga total da memória no contexto ocorre apenas quando o volume é
compacto~\cite{xia2025livesweagent,zhang2025accelopt}.
A granularidade predominante é tarefa (16 estudos, 67\%), alinhada com
a estrutura dos benchmarks em que cada issue equivale a uma tarefa.
Apenas um dos 24 estudos primários reporta compartilhamento de
memória entre agentes~\cite{chen2025sweexp}.
A \autoref{fig:taxonomia} visualiza a distribuição dos estudos no espaço
escopo temporal $\times$ substrato, com cor indicando a política de
controle dominante.

\begin{figure}[htbp]
  \centering
  \resizebox{\textwidth}{!}{% fig:taxonomia — Bubble chart: escopo temporal × substrato × política
% Cada célula mostra a contagem de estudos; cor indica política dominante.
% Requer: tikz (já carregado em main.tex)
\begin{tikzpicture}[
  font=\small,
  cell/.style={
    minimum width=2.0cm, minimum height=1.4cm,
    align=center, inner sep=2pt,
  },
]

% --- Cores de política (aliases de cores definidas em main.tex) ---
\colorlet{polLLM}{archBanco}          % self_control_llm
\colorlet{polHeur}{archHibrido}       % heurística
\colorlet{polApre}{archGit}           % aprendida
\colorlet{polMix}{neutGray}           % célula sem maioria (empate)

% --- Eixos ---
% Colunas: substrato (texto, estruturado, híbrido, executável)
% Linhas: escopo (híbrida, procedural, episódica, semântica)

% Cabeçalhos de coluna
\node[cell, font=\small\bfseries] at (0, 0) {};
\node[cell, font=\small\bfseries] at (2.2, 0) {Texto};
\node[cell, font=\small\bfseries] at (4.4, 0) {Estruturado};
\node[cell, font=\small\bfseries] at (6.6, 0) {Híbrido};
\node[cell, font=\small\bfseries] at (8.8, 0) {Executável};

% Cabeçalhos de linha
\node[cell, font=\small\bfseries, anchor=east] at (-0.2, -1.6) {Híbrida};
\node[cell, font=\small\bfseries, anchor=east] at (-0.2, -3.2) {Procedural};
\node[cell, font=\small\bfseries, anchor=east] at (-0.2, -4.8) {Episódica};
\node[cell, font=\small\bfseries, anchor=east] at (-0.2, -6.4) {Semântica};

% --- Labels de eixo ---
\node[rotate=90, font=\footnotesize\bfseries, anchor=south]
  at (-2.8, -4.0) {Escopo temporal};
\node[font=\footnotesize\bfseries] at (5.5, 0.8)
  {Substrato representacional};

% --- Indicadores de célula vazia (ponto sutil) ---
\foreach \cx/\cy in {
  6.6/-3.2, 8.8/-1.6, 8.8/-4.8, 8.8/-6.4,
  2.2/-6.4, 6.6/-4.8, 6.6/-6.4} {
  \node[circle, fill=black!8, minimum size=0.15cm] at (\cx, \cy) {};
}

% --- Bolhas ---
% Tamanho do círculo proporcional a N; cor = política majoritária na célula

% Híbrida × texto = 5 (Bui, Hosain, Vasilopoulos, Wong, Wu)
% Majoritária LLM (4/5): bui, hosain, wong, wu; heurística: vasilopoulos
\node[circle, fill=polLLM!50, draw=polLLM, minimum size=0.9cm]
  at (2.2, -1.6) {\textbf{5}};

% Híbrida × estruturado = 4 (Chen, Deng, Dong, Zhou)
% Majoritária LLM (3/4): chen, deng, zhou; aprendida: dong
\node[circle, fill=polLLM!50, draw=polLLM, minimum size=0.8cm]
  at (4.4, -1.6) {\textbf{4}};

% Híbrida × híbrido = 5
% (Nadafian, Shen TALM, Tablan, Wang RepoMem, Wang CodeMEM)
% Majoritária LLM (3/5): nadafian, tablan, wang2026codemem;
% heurística: shen2025talm, wang2025repositorymemory
\node[circle, fill=polLLM!50, draw=polLLM, minimum size=0.9cm]
  at (6.6, -1.6) {\textbf{5}};

% Procedural × texto = 1 (Qian) — heurística
\node[circle, fill=polHeur!50, draw=polHeur, minimum size=0.45cm]
  at (2.2, -3.2) {\textbf{1}};

% Procedural × estruturado = 1 (Shen Struct.) — heurística
\node[circle, fill=polHeur!50, draw=polHeur, minimum size=0.45cm]
  at (4.4, -3.2) {\textbf{1}};

% Procedural × executável = 3 (Li GBT, Xia Live-SWE, Zhang DGM)
% Empate: 1 heurística (li), 1 LLM (xia), 1 aprendida (zhang DGM)
\node[circle, fill=polMix!40, draw=polMix, minimum size=0.7cm]
  at (8.8, -3.2) {\textbf{3}};

% Episódica × texto = 3 (Lindenbauer, Su, Zhang AccelOpt) — heurística (3/3)
\node[circle, fill=polHeur!50, draw=polHeur, minimum size=0.7cm]
  at (2.2, -4.8) {\textbf{3}};

% Episódica × estruturado = 1 (Guo EET) — LLM
\node[circle, fill=polLLM!50, draw=polLLM, minimum size=0.45cm]
  at (4.4, -4.8) {\textbf{1}};

% Semântica × estruturado = 1 (Wang MemGovern) — LLM
\node[circle, fill=polLLM!50, draw=polLLM, minimum size=0.45cm]
  at (4.4, -6.4) {\textbf{1}};

% --- Legenda ---
\node[anchor=west] at (0.0, -7.8) {
  \begin{tikzpicture}[font=\footnotesize, baseline=-0.5ex]
    \node[circle, fill=polLLM!50, draw=polLLM,
      minimum size=0.35cm] at (0,0) {};
    \node[anchor=west] at (0.3,0) {LLM controlador};
    \node[circle, fill=polHeur!50, draw=polHeur,
      minimum size=0.35cm] at (3.3,0) {};
    \node[anchor=west] at (3.6,0) {Heurística};
    \node[circle, fill=polMix!40, draw=polMix,
      minimum size=0.35cm] at (5.8,0) {};
    \node[anchor=west] at (6.1,0) {Sem maioria};
  \end{tikzpicture}
};

\end{tikzpicture}
}
  \caption{Distribuição dos 24 estudos primários no espaço
    escopo temporal (linhas) $\times$ substrato representacional
    (colunas).
    O tamanho da bolha é proporcional ao número de estudos em
    cada célula, e a cor indica a política de controle
    majoritária da célula; células sem maioria aparecem em
    cinza.
    A figura evidencia a concentração em bancos de experiência
    com escopo temporal híbrido, célula que reúne o maior número
    de propostas, enquanto substratos como AST-guiado, grafos de
    conhecimento e metáforas git permanecem com adoção isolada.}
  \label{fig:taxonomia}
\end{figure}

\subsection{Métodos de avaliação e benchmarks (QS2)}
\label{sec:qs2}

Dos 28 estudos, 21 utilizam benchmarks, três adotam estudos de caso e
um combina ambos.
A \autoref{tab:benchmarks} lista os benchmarks identificados.
SWE-bench Verified é o benchmark dominante, utilizado por 14 dos 21
estudos com avaliação por benchmark (67\%), o que cria um ponto natural
de comparabilidade entre sistemas.
Em contrapartida, sete estudos empregam benchmarks únicos não
compartilhados por nenhum outro
estudo~\cite{deshpande2025memtrack,dong2026kernelblaster,qian2024iterative,tablan2025smarter,zhang2025accelopt,nadafian2026kapso,wang2026codemem}, o que limita comparações diretas.

\begin{table}[htbp]
  \centering
  \caption{Benchmarks e frequência de uso.}
  \label{tab:benchmarks}
  \begin{tabular}{l r}
    \toprule
    Benchmark & $N$ \\
    \midrule
    SWE-bench Verified    & 14 \\
    SWE-bench Lite        &  2 \\
    SWE-bench Pro         &  2 \\
    WebArena              &  2 \\
    SWE-Bench-CL          &  1 \\
    KernelBench           &  1 \\
    MLE/ALE-Bench         &  1 \\
    DS-1000               &  1 \\
    \midrule
    \multicolumn{2}{l}{\footnotesize
      Outros 4 benchmarks de uso único omitidos.} \\
    \bottomrule
  \end{tabular}
\end{table}

Pass@1 (ou Resolve Rate) é a métrica universal, adotada por 18 dos 21
estudos com benchmark.
Métricas de aprendizado contínuo, como forgetting e backward transfer,
foram propostas por \citeonline{joshi2025swebenchcl} mas ainda não
adotadas por outros estudos.

A comparabilidade direta entre sistemas permanece limitada.
Embora 14 estudos utilizem SWE-bench Verified, as condições
experimentais diferem: pelo menos oito LLMs distintos servem como
backbone, os subsets variam entre 300 e 500 instâncias, e parâmetros
como temperatura e timeout raramente são reportados.
Nenhum estudo compara três ou mais arquiteturas de persistência no mesmo
benchmark com o mesmo modelo, confirmando o Gap~1.

\subsection{Resultados de desempenho (QS3)}
\label{sec:qs3}

Dezenove estudos apresentam comparação controlada contra baseline sem
persistência de memória, todos com ganhos positivos.
A \autoref{tab:desempenho} lista os 17 cujo ganho é mensurável em
pontos percentuais (pp); os dois restantes usam métricas não
comparáveis (\textit{speedup} e score qualitativo).
A magnitude varia de +2,2~pp~\cite{chen2025sweexp} a
+39,0~pp~\cite{li2026traversalaspolicy}, com mediana de
aproximadamente +6,8~pp no SWE-bench Verified.

\begin{table}[htbp]
  \centering
  \caption{Ganhos de desempenho em relação ao baseline sem persistência.
    Ordenados por ganho absoluto.}
  \label{tab:desempenho}
  \begin{tabular}{l l r r r}
    \toprule
    Estudo & Benchmark & Base (\%) & Mem. (\%) & Ganho \\
    \midrule
    Li (GBT) & SWE-b. V. & 34,6 & 73,6 & +39,0~pp \\
    Zhang (DGM) & SWE-b. V. & 20,0 & 50,0 & +30,0~pp \\
    Nadafian (KAPSO) & MLE-Bench & 35,1 & 50,7 & +15,6~pp \\
    Wang (CodeMEM) & CodeIF-B. & $\sim$73 & 85,2 & +12,2~pp \\
    Zhou (ToM-SWE) & Ambig. SWE & 51,9 & 63,4 & +11,5~pp \\
    Su (LbI) & OSWorld & 12,4 & 22,5 & +10,1~pp \\
    Deng (MemCoder) & SWE-b. V. & 68,4 & 77,8 & +9,4~pp \\
    Wang (MemGovern) & SWE-b. V. & 23,2 & 32,6 & +9,4~pp \\
    Su (LbI) & SWE-b. V. & 51,2 & 60,0 & +8,8~pp \\
    Shen (Struct.) & SWE-b. V. & 53,5 & 60,3 & +6,8~pp \\
    Wu (GCC) & SWE-b. V. & 74,0 & 80,2 & +6,2~pp \\
    Wang (RepoMem) & SWE-b. V. & 37,0 & 40,4 & +3,4~pp \\
    Shen (TALM) & BigCodeB. & 50,1 & 53,3 & +3,2~pp \\
    Xia (Live-SWE) & SWE-b. V. & 74,2 & 77,4 & +3,2~pp \\
    Wong (CCA) & SWE-b. Pro & $\sim$56 & 59,0 & +3,0~pp \\
    Guo (EET) & SWE-b. V. & $\sim$32 & $\sim$35 & +2,7~pp \\
    Chen (SWE-Exp) & SWE-b. V. & 70,8 & 73,0 & +2,2~pp \\
    \bottomrule
  \end{tabular}
\end{table}

Os maiores ganhos (acima de +10~pp) ocorrem em sistemas que partem de
baselines mais fracos:
\citeonline{li2026traversalaspolicy} parte de 34,6\%,
\citeonline{zhang2025darwingodel} de 20,0\% e
\citeonline{su2025learnbyinteract} de 12,4\%.
No SWE-bench Verified, baselines acima de 65\% obtêm ganhos entre
+2 e +9~pp, sugerindo rendimentos decrescentes à medida que o
sistema se aproxima do teto do benchmark.

Os ganhos são consistentes entre LLMs.
\citeonline{shen2026structurally} reporta ganho médio de +4,7~pp em
quatro LLMs distintos, e \citeonline{wang2026memgovern} reporta
+4,65~pp em sete LLMs.
Nos estudos com avaliação multi-modelo, a persistência melhora o
desempenho em todos os LLMs testados.
A \autoref{fig:ganhos} visualiza os ganhos por estudo, com cor
indicando o tipo de arquitetura.

\begin{figure}[htbp]
  \centering
  % Source: extracao/dados-extraidos.tsv (study, gain_pp, arch_family)
\begin{tikzpicture}[
  font=\small,
  bar label/.style={anchor=east, font=\footnotesize},
  gain label/.style={anchor=west, font=\footnotesize},
]
\def\barheight{0.35}
\def\scale{0.22}

\def\entries{
  {Li (GBT)/39.0/archOutro},
  {Zhang (DGM)/30.0/archOutro},
  {Nadafian (KAPSO)/15.6/archHibrido},
  {Wang (CodeMEM)/12.2/archAST},
  {Zhou (ToM-SWE)/11.5/archHibrido},
  {Su (LbI) OSWorld/10.1/archBanco},
  {Deng (MemCoder)/9.4/archBanco},
  {Wang (MemGov.)/9.4/archBanco},
  {Su (LbI) SWE-b./8.8/archBanco},
  {Shen (Struct.)/6.8/archBanco},
  {Wu (GCC)/6.2/archGit},
  {Wang (RepoMem)/3.4/archGit},
  {Shen (TALM)/3.2/archBanco},
  {Xia (Live-SWE)/3.2/archOutro},
  {Wong (CCA)/3.0/archHibrido},
  {Guo (EET)/2.7/archBanco},
  {Chen (SWE-Exp)/2.2/archBanco}%
}
\foreach \name/\gain/\col [count=\i from 0] in \entries {
  \pgfmathsetmacro{\y}{-\i * 0.65}
  \pgfmathsetmacro{\w}{\gain * \scale}
  \node[bar label] at (-0.2, \y) {\name};
  \fill[\col!80] (0, \y - \barheight/2)
    rectangle (\w, \y + \barheight/2);
  \draw[\col!60!black, line width=0.3pt] (0, \y - \barheight/2)
    rectangle (\w, \y + \barheight/2);
  \node[gain label] at (\w + 0.1, \y)
    {+\pgfmathprintnumber[fixed, precision=1]{\gain}};
}
\pgfmathsetmacro{\ybot}{-16 * 0.65 - 0.5}
\draw[thick, -Stealth] (0, 0.5) -- (0, \ybot);
\draw[thick] (0, \ybot) -- (9.0, \ybot);
\foreach \v in {0, 5, 10, 15, 20, 25, 30, 35, 40} {
  \pgfmathsetmacro{\x}{\v * \scale}
  \draw (\x, \ybot) -- (\x, \ybot - 0.15);
  \node[font=\scriptsize, below] at (\x, \ybot - 0.15) {\v};
}
\node[font=\footnotesize, below] at (4.4, \ybot - 0.7)
  {Gain (percentage points)};
\pgfmathsetmacro{\yleg}{\ybot - 1.8}
\fill[archBanco!80] (0.5, \yleg) rectangle (0.85, \yleg+0.25);
\node[anchor=west, font=\footnotesize] at (0.95, \yleg+0.125)
  {Experience bank};
\fill[archHibrido!80] (3.2, \yleg) rectangle (3.55, \yleg+0.25);
\node[anchor=west, font=\footnotesize] at (3.65, \yleg+0.125)
  {Hybrid};
\fill[archGit!80] (5.6, \yleg) rectangle (5.95, \yleg+0.25);
\node[anchor=west, font=\footnotesize] at (6.05, \yleg+0.125)
  {Git metaphor};
\pgfmathsetmacro{\ylegb}{\yleg - 0.5}
\fill[archAST!80] (0.5, \ylegb) rectangle (0.85, \ylegb+0.25);
\node[anchor=west, font=\footnotesize] at (0.95, \ylegb+0.125)
  {AST-guided};
\fill[archOutro!80] (3.2, \ylegb) rectangle (3.55, \ylegb+0.25);
\node[anchor=west, font=\footnotesize] at (3.65, \ylegb+0.125)
  {Other};

\end{tikzpicture}

  \caption{Ganho em pontos percentuais ($\Delta$~pp) sobre o
    baseline sem persistência, para os 17 estudos com métricas
    comparáveis.
    Cada barra corresponde a um estudo, com cor indicando o tipo
    de arquitetura de memória.
    A mediana cross-study é de 6,8~pp; os maiores ganhos (acima
    de 10~pp) concentram-se em sistemas que partem de baselines
    abaixo de 35\%, enquanto baselines acima de 65\% apresentam
    ganhos entre 2 e 9~pp, padrão consistente com rendimentos
    decrescentes.
    Os estudos Dong (KernelBlaster) e Tablan (Spark) são
    omitidos por usarem métricas não comparáveis (speedup e
    score qualitativo, respectivamente).}
  \label{fig:ganhos}
\end{figure}

Cinco estudos não apresentam comparação
controlada~\cite{bui2026opendev,vasilopoulos2026codified,qian2024iterative,hosain2025xolver,zhang2025accelopt}, por razões
como ausência de benchmark formal, natureza observacional do estudo
de caso ou baseline incomparável.

\subsection{Custo e eficiência (QS4)}
\label{sec:qs4}

O reporte de dados de custo é heterogêneo.
Apenas sete estudos fornecem breakdown completo com tokens, custo
monetário e comparação com baseline.
Outros nove reportam dados parciais (custo ou tokens, sem ambos).
Os 12 estudos restantes (43\%) não disponibilizam nenhum dado de custo,
confirmando o Gap~2.

A \autoref{tab:custo} apresenta o custo monetário por tarefa nos
estudos que disponibilizam essa informação.
A variação é ampla: de US\$\,0,01 por
tarefa~\cite{guo2026eet} a US\$\,7,27~\cite{wang2026memgovern},
refletindo diferenças de LLM, benchmark e arquitetura.

\begin{table}[htbp]
  \centering
  \caption{Custo monetário por tarefa nos estudos que reportam
    dados monetários, agrupados por configuração de LLM quando o
    estudo apresenta múltiplas. Os sete estudos com detalhamento
    completo de custo (Chen, Guo, MemGovern, Wu, Xia, Zhang
    AccelOpt e Zhou) aparecem com pelo menos uma linha; Wang
    RepoMem é incluído apesar do detalhamento parcial por
    apresentar custo monetário por tarefa.}
  \label{tab:custo}
  \begin{tabular}{l r l l}
    \toprule
    Estudo & Custo/tarefa & LLM & Benchmark \\
    \midrule
    Guo (EET) & US\$\,0,01 & GPT-5-mini & SWE-b. V. \\
    Zhou (ToM-SWE) & US\$\,0,02 & GPT-5-nano & Stateful SWE \\
    Xia (Live-SWE) & US\$\,0,05--0,73 & vários & SWE-b. V./Pro \\
    Wang (MemGov., Qwen) & US\$\,0,07 & Qwen3-30B & SWE-b. V. \\
    Chen (SWE-Exp) & US\$\,0,13 & DeepSeek-V3 & SWE-b. V. \\
    Wu (GCC) & US\$\,0,57--1,21 & vários & SWE-b. Lite \\
    Wang (RepoMem) & US\$\,0,58--0,89 & n/d & SWE-b. V. \\
    Zhang (AccelOpt) & US\$\,6,87--15,95 & open-source & NKIBench \\
    Wang (MemGov., Claude) & US\$\,7,27 & Claude 4 S. & SWE-b. V. \\
    \midrule
    \multicolumn{4}{l}{\footnotesize n/d: não declarado no paper original.} \\
    \bottomrule
  \end{tabular}
\end{table}

A relação entre memória e custo segue dois padrões distintos.
Sistemas que usam memória para substituir contexto redundante reduzem
custos: \citeonline{guo2026eet} obtém $-31{,}8$\% no custo total,
\citeonline{li2026traversalaspolicy} reduz tokens em 39\% e
\citeonline{shen2025talm} consome aproximadamente metade dos tokens de
seu baseline no ClassEval.
Em contrapartida, sistemas que adicionam contexto via retrieval
aumentam custos marginalmente:
\citeonline{chen2025sweexp} reporta $+7{,}5$\% em tokens e
\citeonline{wang2026memgovern} $+4{,}8$\% com Claude 4 Sonnet.
Os ganhos de desempenho nesses casos são proporcionais ao custo
adicional.
A \autoref{fig:custo} ilustra essa dualidade.

\begin{figure}[htbp]
  \centering
  \resizebox{\textwidth}{!}{% Requer: tikz (já carregado em main.tex)
% Source: extracao/dados-extraidos.tsv (study, cost_impact, cost_mechanism)
\begin{tikzpicture}[
  font=\small,
  system box/.style={
    draw=#1!70, fill=#1!10, rounded corners=3pt,
    minimum width=3.8cm, minimum height=0.7cm,
    align=center, font=\footnotesize,
  },
  system box/.default=archBanco,
  group label/.style={
    font=\small\bfseries, align=center,
  },
  arrow down/.style={-Stealth, thick, black!60},
]
\node[group label, text=posGain] at (-3.5, 0)
  {Substituição\\[-2pt]{\footnotesize (reduz custo)}};
\node[group label, text=negLoss] at (3.5, 0)
  {Adição\\[-2pt]{\footnotesize (aumenta custo)}};

\draw[dashed, black!30, thick] (0, 0.5) -- (0, -7.5);
\node[system box=posGain] at (-3.5, -1.0)
  {EET: $-31{,}8$\% custo};
\node[system box=posGain] at (-3.5, -2.0)
  {GBT: $-39$\% tokens};
\node[system box=posGain] at (-3.5, -3.0)
  {TALM: $\sim$$-50$\% tokens};
\node[system box=posGain] at (-3.5, -4.0)
  {KernelBlaster: $2{,}4\times$ menos};
\node[system box=posGain] at (-3.5, -5.0)
  {CodeMEM: $-30$k tok./round};
\node[system box=posGain] at (-3.5, -6.0)
  {CCA: $-11$k tok./tarefa};
\node[system box=posGain] at (-3.5, -7.0)
  {AccelOpt: $26\times$ mais barato};
\node[system box=negLoss] at (3.5, -2.5)
  {SWE-Exp: $+7{,}5$\% tokens};
\node[system box=negLoss] at (3.5, -3.5)
  {MemGovern: $+4{,}8$\% custo};
\node[system box=negLoss] at (3.5, -4.5)
  {Live-SWE: $+4$--$21$\%};
\node[system box=negLoss] at (3.5, -5.5)
  {RepoMem: custo $\uparrow$ em tarefas difíceis};
\node[font=\scriptsize, text=black!60, align=center,
  text width=3.5cm] at (-3.5, -8.2)
  {Memória substitui contexto\\redundante ou evita\\retrabalho};
\node[font=\scriptsize, text=black!60, align=center,
  text width=3.5cm] at (3.5, -7.2)
  {Memória adiciona contexto\\via retrieval; ganho de\\desempenho proporcional};

\end{tikzpicture}
}
  \caption{Dois padrões opostos de impacto da memória no custo
    operacional, observados nos estudos com dados de custo
    completos.
    À esquerda, sistemas que substituem contexto redundante por
    memória condensada (EET, Traversal-as-Policy, TALM) reduzem
    o custo total entre 32\% e 50\%.
    À direita, sistemas que adicionam contexto via retrieval
    (SWE-Exp, MemGovern com Claude~4~Sonnet) elevam o custo
    entre 5\% e 21\%.
    Cada barra representa a variação percentual de tokens ou
    custo monetário em relação ao baseline sem memória do
    próprio estudo.}
  \label{fig:custo}
\end{figure}

\citeonline{guo2026eet} apresenta o melhor equilíbrio custo-desempenho
entre os estudos com dados disponíveis:
$-31{,}8$\% de custo com perda máxima de $-0{,}2$~pp, ao identificar
oportunidades de terminação antecipada em 11,3\% das issues.
\citeonline{zhang2025accelopt} é 26 vezes mais barato que Claude
Sonnet~4 com desempenho equivalente em seu domínio.
Nenhum estudo constrói fronteira de Pareto explícita entre acurácia e
custo.
Os dados disponíveis indicam que alguns
sistemas~\cite{guo2026eet,li2026traversalaspolicy} superam outros
simultaneamente em custo e desempenho, sugerindo que a fronteira
eficiente atual ainda pode ser melhorada.

\subsection{Complexidade versus eficácia (QS5)}
\label{sec:qs5}

Nove dos 28 estudos (32\%) reportam cenários onde a memória degrada
o desempenho.
A \autoref{tab:degradacao} sumariza esses casos.

\begin{table}[htbp]
  \centering
  \caption{Cenários de degradação por memória.}
  \label{tab:degradacao}
  \begin{tabularx}{\textwidth}{p{2.8cm} @{\hspace{1em}} p{2.8cm} @{\hspace{1em}} X}
    \toprule
    Estudo & Mecanismo & Cenário \\
    \midrule
    Deshpande & ruído & Mem0 e Zep mantêm redundância de
      \textit{tool calls} em $\sim$21\%, sem melhoria sobre baseline \\
    Joshi & infraestrutura & Harness incompatível com mecanismo de
      memória \\
    Lindenbauer & ruído & Itens ruidosos desviam exploração
      ($-11$~pp) \\
    Chen & ruído & $k > 1$ experiências degrada (42,0\% para
      39,0\%) \\
    Shen (TALM) & overhead & Retrieval não compensa em tarefas
      simples \\
    Vasilopoulos & manutenção & Knowledge-to-code ratio de 24\% requer
      1--2h/semana \\
    Xia & overhead & LLMs fracos degradam $-68{,}2$\% com
      \textit{tool creation} \\
    Srivastava & envenenamento & 9,1\% de experiências envenenadas
      comprometem 47,9\% dos retrievals \\
    Zhou & viés & User simulator introduz over-compliance \\
    \bottomrule
  \end{tabularx}
\end{table}

Quatro mecanismos de degradação emergem dos dados: ruído na recuperação,
onde experiências irrelevantes ou excessivas desviam o
agente~\cite{lindenbauer2025ctimrover,chen2025sweexp,deshpande2025memtrack}; overhead desnecessário, onde o custo do
retrieval não compensa em tarefas simples ou com modelos
fracos~\cite{shen2025talm,xia2025livesweagent}; envenenamento da
memória, uma vulnerabilidade de segurança em stores
abertos~\cite{srivastava2025memorygraft}; e incompatibilidade entre o harness de avaliação e a ordenação
cronológica de tarefas~\cite{joshi2025swebenchcl}, uma limitação
experimental que o próprio estudo reconhece.

Não há evidência de que maior complexidade arquitetural produza ganhos
sistematicamente maiores.
\citeonline{zhang2025accelopt}, com fila simples de experiências, e
\citeonline{guo2026eet}, com heurística de terminação, obtêm ganhos
expressivos em seus respectivos domínios sem recorrer a arquiteturas
multi-tier como KAPSO~\cite{nadafian2026kapso} e
GCC~\cite{wu2025gcc}.
O fator com maior associação aparente ao ganho é a distância entre o
baseline e o
potencial do benchmark: memória simples aplicada a baseline
fraco~\cite{zhang2025darwingodel} produz +30~pp, enquanto memória
sofisticada em baseline forte~\cite{chen2025sweexp} resulta em
+2,2~pp.

A calibração do retrieval é determinante.
\citeonline{chen2025sweexp} reporta que recuperar mais de uma
experiência degrada o resultado (42,0\% para 39,0\% com $k = 4$), e
\citeonline{shen2026structurally} mostra que isolamento por categoria supera
retrieval global por +2,3~pp.

\subsection{Análises de sensibilidade e certeza}
\label{sec:sensibilidade}

Duas análises de sensibilidade verificaram a estabilidade dos achados.
A primeira (S1) reteve apenas os cinco estudos publicados em venues com
revisão por pares.
Os achados principais (ganhos positivos, rendimentos decrescentes,
degradação por ruído) se sustentam qualitativamente, porém a base
empírica de QS4 cai de sete para dois estudos ao excluir
EET~\cite{guo2026eet}, MemGovern~\cite{wang2026memgovern} e
GCC~\cite{wu2025gcc}.
A segunda (S2) excluiu os dois estudos com pontuação inferior a 3,0;
a remoção não alterou nenhum achado quantitativo, pois ambos contribuem
apenas descrições qualitativas.

A \autoref{tab:certeza} apresenta a avaliação de certeza por achado.
Cada linha recebe também uma classificação de tipo de evidência:
\textbf{A} para dados diretos extraídos do TSV ou JSON (contagens,
percentuais, verificação exaustiva), \textbf{B} para padrões
cross-study agregados a partir de múltiplos estudos e \textbf{C}
para caracterizações de mecanismo ou inferências causais cuja base
empírica é inerentemente interpretativa.
Dos 19 achados codificados, 13 recebem certeza alta, cinco recebem
certeza moderada e um recebe certeza baixa; em paralelo, 13 são
classificados no tipo A, quatro no tipo B e dois no tipo C.
O único achado de certeza baixa é QS1-D (substratos executáveis,
baseado em apenas três estudos), enquanto QS5-B (relação
complexidade versus ganho) recebe certeza moderada por depender
de comparação entre estudos com benchmarks distintos.

\begin{table}[htbp]
  \centering
  \caption{Certeza e tipo de evidência por achado. Tipo A: dados
  diretos (contagens, percentuais). Tipo B: padrão cross-study
  agregado. Tipo C: caracterização de mecanismo ou inferência causal.}
  \label{tab:certeza}
  \begin{tabularx}{\textwidth}{l c l X}
    \toprule
    Achado & Tipo & Certeza & Justificativa \\
    \midrule
    QS1-A: bancos dominam (46\%) & A & alta & 11 estudos, contagem direta \\
    QS1-B: escopo híbrido (58\%) & A & alta & 14 estudos \\
    QS1-C: LLM controla memória (54\%) & A & alta & 13 estudos,
      classificação direta \\
    QS1-E: recuperação híbrida (46\%) & A & alta & 11 estudos \\
    QS1-F: compartilhamento inexplorado & A & alta & 23/24 reportam
      ausência \\
    QS2-A: SWE-bench V. dominante & A & alta & 14/21 estudos \\
    QS2-D: Gap~1 confirmado & A & alta & verificação exaustiva \\
    QS3-A: ganhos universais & B & alta & 19 comparações, zero
      inversões \\
    QS4-A: 43\% sem custo & A & alta & contagem direta \\
    QS4-D: ausência de fronteira de Pareto & A & alta & verificação
      exaustiva \\
    QS5-A: 32\% degradação & A & alta & 9 estudos com evidência \\
    T-A: campo dominado por preprints (82\%) & A & alta & contagem
      direta \\
    T-C: Gap~1 permanece confirmado & A & alta & verificação exaustiva \\
    \midrule
    QS3-B: rendimentos decrescentes & B & moderada & padrão observado,
      não testado formalmente \\
    QS3-C: consistência cross-LLM & B & moderada & 2 estudos
      multi-LLM \\
    QS4-B: dois padrões de custo & B & moderada & 12 estudos parciais \\
    QS5-B: complexidade $\neq$ ganho & C & moderada & inferência
      cross-study, benchmarks heterogêneos \\
    QS5-D: calibração retrieval & C & moderada & 2 estudos diretos \\
    \midrule
    QS1-D: substratos executáveis & A & baixa & 3 estudos apenas \\
    \bottomrule
  \end{tabularx}
\end{table}

A predominância de preprints no corpus (82\%) introduz risco de viés de
reporte, particularmente para QS4 (dados de custo, frequentemente
ausentes em preprints) e QS2 (seleção de benchmarks, menos padronizada
em trabalhos não revisados por pares).
A análise de sensibilidade S1, restrita aos cinco estudos publicados em
venues com revisão por pares, não inverteu nenhum achado, o que sugere
que o risco de viés está contido para os achados de alta certeza.
Para os achados de certeza moderada e baixa (\autoref{tab:certeza}),
a influência de preprints não pode ser descartada.

\Needspace{14\baselineskip}
\section{Discussion}
\label{sec:discussao}

% PRISMA Item 17: Discussion of results in the context of other evidence
% PRISMA Item 18: Limitations of the review

\subsection{Main findings}
\label{sec:principais-achados}

The 28 mapped studies reveal a field concentrated around one dominant
architecture: experience banks that distil resolution trajectories into
reusable entries (46\% of primary studies,
\autoref{tab:arquiteturas}), with the LLM itself acting as the memory
controller (54\%) and evaluation centred on SWE-bench Verified (67\% of
studies with a benchmark, \autoref{tab:benchmarks}).
This convergence, however, does not imply maturity.
The field remains dominated by preprints (82\%), and the 28 studies use
at least eight distinct LLMs as backbones.
Controlled comparisons across architectures on the same benchmark with
the same model are absent from the corpus, so the observed convergence
reflects independent author choices rather than an empirical validation
of experience banks over alternatives such as knowledge graphs or Git
metaphors.

All 19 studies with controlled comparison report improvements over the
no-memory baseline (\autoref{tab:desempenho}), with a median gain of
6.8~pp on SWE-bench Verified and an amplitude of +2.2~pp to +39.0~pp.
The data indicate a diminishing-returns pattern:
on SWE-bench Verified, systems with baselines below 35\% obtain gains
above
10~pp~\citep{li2026traversalaspolicy,zhang2025darwingodel},
while baselines above 65\% obtain gains of 2 to
9~pp~\citep{deng2026memcoder,wu2025gcc,chen2025sweexp}.
The cost impact is bimodal: systems that replace context with memory
cut tokens by up to
39\%~\citep{li2026traversalaspolicy}.
Systems that add context via retrieval raise cost by +5 to
+21\%~\citep{chen2025sweexp,wang2026memgovern,xia2025livesweagent}.

The risk of degradation is not negligible.
Nine of the 28 studies (32\%) report scenarios where memory hurts
performance (\autoref{tab:degradacao}).
Retrieval noise is the mechanism cited in at least three of the nine
studies with degradation:
\citet{lindenbauer2025ctimrover} document a drop of 11~pp from unfiltered
item injection into context,
\citet{chen2025sweexp} document degradation with $k > 1$ retrieved
experiences, and
\citet{deshpande2025memtrack} identify that commercial memory backends
(Mem0, Zep) do not improve performance and keep tool-call redundancy at
roughly 21\%, with no reduction relative to the no-memory baseline.
The poisoning vulnerability documented by
\citet{srivastava2025memorygraft}, in which 9.1\% of malicious
experiences compromise 47.9\% of retrievals under dual-channel
retrieval (lexical + vector), indicates that memory stores without
provenance checks can constitute an attack surface.

\subsection{Implications for research}
\label{sec:implicacoes-pesquisa}

The central gap identified by this review is the absence of head-to-head
comparisons.
None of the 28 studies evaluates three or more persistence architectures
on the same benchmark, with the same model, and under the same compute
budget.
Each system reports gains against its own no-memory baseline, but
experimental conditions differ in model, framework, instance subset, and
number of attempts.
This fragmentation prevents any conclusion on which architecture is
superior, because the observed gains may reflect implementation
differences as much as differences in memory architecture.
A comparative study evaluating at least three representative
architectures under controlled conditions would fill this gap.

Cost underreporting likewise limits the practical utility of the
results.
Only seven studies provide a complete
breakdown~\citep{chen2025sweexp,guo2026eet,wang2026memgovern,wu2025gcc,xia2025livesweagent,zhang2025accelopt,zhou2025tomswe},
and 12 of the 28 (43\%) report no cost data at all.
Without standardised cost-per-task metrics (tokens consumed, monetary
cost, wall-clock time), practitioners cannot assess the trade-off of
adopting persistence.
Constructing a Pareto frontier between accuracy and cost, absent in
every analysed study, is a research direction with immediate applied
value.

Evidence on the relationship between architectural complexity and gain
is partial.
The data indicate that simple systems, such as
AccelOpt~\citep{zhang2025accelopt} with an experience queue and
EET~\citep{guo2026eet} with heuristic early termination, obtain
relevant gains in their respective domains without resorting to
multi-tier architectures.
However, this comparison is between distinct studies: benchmarks and
models differ across systems.
Only a study controlling the complexity variable internally could
confirm whether architectural simplicity is indeed sufficient or
whether scenarios exist where sophistication is justified.

Three studies propose memory as a directly executable artefact, an
emerging trend with still-limited evidence:
behaviour trees~\citep{li2026traversalaspolicy},
scripts created at run time~\citep{xia2025livesweagent}, and
whole agents~\citep{zhang2025darwingodel}.
The first two report substantial gains (+39.0~pp and +3.2~pp,
respectively), but the certainty of this finding is low given the
small number of studies and the absence of independent replication.

\subsection{Implications for practice}
\label{sec:implicacoes-pratica}

For teams that wish to incorporate knowledge persistence into coding
agents, the data in this review suggest a conservative starting point:
experience banks with similarity-based retrieval.
This is the dominant architecture (46\%), has reference open-source
implementations, and yields gains of 2 to 9~pp in every study with
baselines above 65\%.
More complex systems (knowledge graphs, hierarchical memory) do not
present, in the available data, clear evidence of systematically larger
gains, though this comparison is limited by benchmark and model
differences across studies.

The available data suggest that retrieval-mechanism calibration may
matter more than architectural sophistication.
\citet{chen2025sweexp} report that retrieving more than one experience
($k > 1$) degrades performance from 42.0\% to 39.0\%, and
\citet{shen2026structurally} report that category isolation outperforms
global retrieval by 2.3~pp.
Both results point in the same direction: selectivity in retrieval
influences the effective gain, and excess retrieved context can
introduce noise.

For already-deployed systems, experience-guided early termination
offers cost reduction with minimal risk.
\citet{guo2026eet} report 31.8\% savings in total cost with a maximum
loss of 0.2~pp, by identifying termination opportunities in 11.3\% of
instances.
This approach is agent-agnostic and requires no changes to the
underlying memory architecture.

Of the 24 primary studies, 17 share a complete repository and three
offer partial access, totalling 83\% with some degree of
reproducibility.

\subsection{Limitations of this review}
\label{sec:limitacoes}

The limitations of this review are organised into four categories of
validity threat: internal validity (rigour of the selection and
extraction processes), external validity (generalisation of the
findings beyond the corpus), construct validity (operationalisation of
the measured concepts), and conclusion validity (stability of
cross-study inferences).

Screening was conducted by a single researcher, which is an
acknowledged threat in systematic reviews and precludes the computation
of formal agreement indices such as Cohen's $\kappa$.
This limitation was mitigated by eligibility criteria calibrated in a
pilot study and refined through a stress test against five comparable
systematic reviews in
LLM/SE~\citep{hou2024llmse,zhang2024unifying,jin2024llmagents,wang2024llmagents,liu2024llmagents},
a process that produced 17 documented borderline cases and a decision
flow chart with eight sequential questions.
Each screening decision was recorded in a versioned spreadsheet with
per-criterion rationale, permitting retrospective audit of the
exclusions.

Data extraction was performed by the same researcher, introducing risk
of interpretation bias in qualitative fields such as architecture-type
classification and degradation mechanism.
Four mitigations reduced this risk: (i) programmatic validation of 448
enum values (16 fields $\times$ 28 studies), with error rate below 1\%
after corrections; (ii) spot-check of 40 interpretive fields across 10
selected studies, as detailed in \autoref{sec:extracao}; (iii) factual
fidelity digests generated from the full text of each study, used as
reference when drafting the manuscript's interpretive claims; and (iv)
re-verification of the \textit{negative\_result} fields against the
original PDF.
Even with these mitigations, formal reproducibility of the
classification in interpretive fields was not assessed.

The preprint prevalence (82\% of the corpus) reduces confidence in this
review's findings.
The field of knowledge persistence for coding agents is recent, with
publications concentrated between 2024 and 2026, and only five studies
have been published in peer-reviewed venues: two at
conferences~\citep{wang2026improving,zhang2025darwingodel} and three at
workshops~\citep{deshpande2025memtrack,shen2025talm,lindenbauer2025ctimrover}.
Sensitivity analysis S1 (\autoref{sec:sensibilidade}) shows that,
when preprints are excluded, no qualitative conclusion inverts, but the
evidence base for cost (SQ4) drops from seven to two studies and the
detailed taxonomy loses empirical support.

Searches were limited to two databases (Scopus and arXiv), complemented
by citation chasing.
Studies published exclusively in uncovered bases (for example, IEEE
Xplore, ACM Digital Library) may have been missed.
The rationale for this choice, based on \citet{mourao2020hybrid}, is
documented in \autoref{sec:fontes}.

The 14 studies that use SWE-bench Verified differ in base model, agent
framework, instance subset, and run conditions (\autoref{sec:qs2}).
The construct ``agent performance'' is not operationalised uniformly
across studies, which limits direct comparability of resolution
percentages reported as if they were equivalent quantities.
The taxonomy proposed in this review for persistence architectures also
embeds interpretive choices: families such as ``hybrid'' group
heterogeneous mechanisms whose fine differentiation would require
metadata absent from published texts.

Heterogeneity across studies precludes formal meta-analysis and limits
inferences to qualitative patterns synthesised narratively.
Classifying findings into evidence types A, B, and C
(\autoref{tab:certeza}) makes the empirical support of each inference
explicit: type A for direct counts, type B for aggregated cross-study
patterns, and type C for mechanism characterisations.
The two sensitivity analyses conducted (S1, preprint exclusion;
S2, exclusion of studies scoring below 3.0) did not invert any
quantitative finding, reinforcing the stability of the synthesis under
alternative inclusion criteria.

\subsection{Research agenda}
\label{sec:agenda-pesquisa}

The three gaps confirmed by this review define the field's immediate
agenda.
Gap~1 --- the absence of head-to-head comparisons --- calls for an
experiment that evaluates at least three representative architectures
on the same benchmark, with the same model, and under the same compute
budget.
Gap~2 --- cost underreporting --- requires that such an experiment
include an integrated cost-per-task analysis, producing an explicit
Pareto frontier between accuracy and cost.

Metric standardisation is a prerequisite for future comparability.
The continual-learning metrics proposed by
\citet{joshi2025swebenchcl} (\textit{CL-Score},
\textit{forward/backward transfer}, \textit{forgetting}) have been
adopted by no other study in the corpus.
Adopting standardised cost metrics (tokens per task, monetary cost,
wall-clock time) and knowledge-retention metrics would enable
cross-study comparisons that are currently unfeasible.

Only one of the 24 primary studies reports memory sharing across
agents~\citep{chen2025sweexp}, and even there it is limited to
experience reuse within the same framework.
Sharing protocols across heterogeneous agents remain without empirical
evaluation.
The vulnerabilities flagged by \citet{srivastava2025memorygraft}
indicate that any such protocol must address the risk of poisoning in
open stores before it can be recommended for production.

\section{Conclusão}
\label{sec:conclusao}

Esta revisão sistemática mapeou 28 estudos sobre persistência de
conhecimento em agentes de codificação baseados em IA, todos
publicados a partir de 2024.
O campo é recente e predominantemente composto por preprints (82\%),
mas já apresenta convergência em torno de padrões arquiteturais e
de avaliação recorrentes.
A persistência de memória produz ganhos em todos os cenários
testados, porém com rendimentos decrescentes, riscos de degradação
documentados e lacunas persistentes na documentação de custos.

As cinco questões secundárias receberam respostas com diferentes
graus de certeza.
Quanto à taxonomia (QS1), bancos de experiência predominam (46\%),
o LLM atua como controlador em 54\% dos sistemas, e o escopo
temporal híbrido é o mais frequente (58\%).
Sobre avaliação (QS2), SWE-bench Verified concentra 67\% dos
estudos com benchmark, porém a comparabilidade é limitada por
diferenças de modelo, framework e condições experimentais.
Os ganhos de desempenho (QS3) são universalmente positivos nas 19
comparações controladas, com mediana de 6,8~pp e amplitude de
+2,2 a +39,0~pp, e rendimentos decrescentes em baselines acima
de 65\%.
Em relação ao custo (QS4), 43\% dos estudos não reportam dados de
custo, e o impacto da memória é bimodal: sistemas que substituem
contexto reduzem tokens, enquanto sistemas que adicionam contexto
elevam o custo em aproximadamente 5 a 21\%.
Sobre complexidade e eficácia (QS5), 32\% dos estudos reportam
degradação por memória, com ruído na recuperação como mecanismo
principal, e não há evidência de que arquiteturas mais complexas
produzam ganhos sistematicamente superiores.

Nenhum dos 28 estudos compara arquiteturas de persistência
diretamente.
Essa ausência, somada à subnotificação de custos e à incerteza
sobre o papel da complexidade arquitetural, impede decisões
informadas sobre qual abordagem adotar em cada contexto.
A segunda etapa desta dissertação endereçará a primeira dessas
lacunas por meio de um experimento comparativo com ao menos três
arquiteturas representativas, sob o mesmo benchmark, o mesmo
modelo e com análise integrada de custo por tarefa.


\appendix
\section{Validação da estratégia de busca}
\label{ap:validacao-busca}

Este apêndice documenta o processo de validação aplicado à estratégia de busca antes da execução das consultas finais nas bases.
A validação combina dois procedimentos: auditoria das strings contra revisões sistemáticas comparáveis e verificação de recall por meio de um conjunto de artigos-semente operando como Quasi-Gold Standard (QGS).
A última subseção apresenta o reporte completo do PRISMA-S~\cite{rethlefsen2021prismas}.

\subsection{Auditoria pré-execução}
\label{ap:auditoria}

A auditoria pré-execução consultou seis revisões sistemáticas e mapeamentos publicados entre 2024 e 2026 sobre agentes de codificação baseados em \textit{Large Language Models} (LLM), agentes de software autônomos e técnicas de memória em agentes de IA.
Para cada revisão foram extraídos: as strings completas (quando disponíveis), o conjunto de termos do bloco de agentes e do bloco de memória, e os sistemas representativos catalogados.
Em paralelo, foi compilada uma lista de 16 sistemas conhecidos da literatura corrente como casos-teste de recall, e dez formulações alternativas de termos foram testadas individualmente em consultas exploratórias no Scopus e no arXiv.
A \autoref{tab:auditoria-slrs} resume as revisões consultadas e a contribuição de cada uma para o refinamento dos blocos.

\begin{table}[htbp]
\centering
\caption{Revisões sistemáticas auditadas para calibração das strings de busca.}
\label{tab:auditoria-slrs}
\small
\begin{tabular}{p{3.2cm}p{1cm}p{2.2cm}p{1cm}p{4.5cm}}
\toprule
\textbf{Revisão} & \textbf{Ano} & \textbf{Venue} & \textbf{N} & \textbf{Contribuição ao refinamento} \\
\midrule
Hou et al.~\cite{hou2024llmse} & 2024 & TOSEM & 395 & Termos de bloco \textit{coding}/\textit{software engineering}; cobertura de venues. \\
Liu et al.~\cite{liu2024llmagents} & 2024 & TOSEM & 117 & Vocabulário de agentes de software autônomos. \\
Zhang et al.~\cite{zhang2024unifying} & 2024 & SCIS & 52 & Termos de planejamento e memória em agentes LLM. \\
Jin et al.~\cite{jin2024llmagents} & 2024 & arXiv & 90 & Variantes de \textit{LLM-based agent}; sinônimos de bloco C1. \\
Wang et al.~\cite{wang2024llmagents} & 2024 & ASE & 106 & Confirmação do termo \textit{language agent}; vocabulário de papéis. \\
Du et al.~\cite{du2026memory} & 2026 & arXiv & --- & Taxonomia de memória em agentes de IA; vocabulário de bloco C2. \\
\bottomrule
\end{tabular}
\end{table}

A auditoria identificou três lacunas no rascunho inicial das strings: ausência de variantes sem sufixo \textit{-ing} (\textit{code assistant}), ausência do termo emergente \textit{agentic coding} (já adotado como categoria de venue em 2025) e ausência da forma \textit{SWE agent} como nome próprio de classe.
As lacunas correspondentes no bloco de memória abrangiam \textit{persistent memory}, \textit{persistent context} e \textit{cross-task}, recorrentes nos papers-semente e nas revisões auditadas.

\subsection{Recuperação de artigos-semente (Quasi-Gold Standard)}
\label{ap:qgs}

A validação por Quasi-Gold Standard consiste em verificar se uma estratégia de busca recupera um conjunto pré-selecionado de artigos representativos do escopo, identificados independentemente da estratégia sendo validada~\cite{mourao2020hybrid}.
Foram definidos sete artigos-semente, escolhidos por cobrirem diferentes arquiteturas de persistência (banco de experiência, memória episódica, memória orientada a commits, contexto inspirado em \textit{git}) e por incluírem ao menos um resultado negativo.
A estratégia passou por dois ciclos de refinamento documentados no
protocolo versionado, elevando a recuperação de quatro para seis dos
sete artigos-semente; o caso restante (CodeMEM) foi resgatado pelo rastreamento de citações, conforme descrito na seção de método.
A \autoref{tab:qgs-recuperacao} consolida o status de recuperação por base.

\begin{table}[htbp]
\centering
\caption{Recuperação dos artigos-semente por base e método. \textit{Cit.} indica resgate via rastreamento de citações.}
\label{tab:qgs-recuperacao}
\small
\begin{tabular}{p{3.5cm}p{1cm}p{1.5cm}p{1.2cm}p{1.5cm}p{2cm}}
\toprule
\textbf{Artigo-semente} & \textbf{Ano} & \textbf{Scopus} & \textbf{arXiv} & \textbf{S2 ID} & \textbf{Recuperado} \\
\midrule
Du (survey de memória)~\cite{du2026memory} & 2026 & não & sim & sim & Sim (string) \\
Joshi et al.\ (SWE-Bench-CL)~\cite{joshi2025swebenchcl} & 2025 & não & sim & sim & Sim (string) \\
Wang et al.\ (CodeMEM)~\cite{wang2026codemem} & 2026 & não & não & sim & Sim (Cit.) \\
Deng et al.\ (MemCoder)~\cite{deng2026memcoder} & 2026 & não & sim & sim & Sim (string) \\
Wu et al.\ (GCC)~\cite{wu2025gcc} & 2025 & não & sim & sim & Sim (string) \\
Wang et al.\ (Repository Memory)~\cite{wang2026improving} & 2026 & não & sim & sim & Sim (string) \\
Lindenbauer et al.\ (CTIM-Rover)~\cite{lindenbauer2025ctimrover} & 2025 & não & sim & sim & Sim (string) \\
\bottomrule
\end{tabular}
\end{table}

A ausência total dos artigos-semente no Scopus reflete uma característica do campo: a maioria dos sistemas analisados é composta de \textit{preprints} de 2025--2026 ainda fora do ciclo de indexação de venues tradicionais.
Esse diagnóstico foi confirmado por consultas auxiliares no Scopus (busca por nomes próprios de sistemas), que retornaram zero resultados, e justifica a centralidade do rastreamento de citações como terceiro método PRISMA na estratégia adotada.

\subsection{Conformidade PRISMA-S}
\label{ap:prisma-s}

A \autoref{tab:prisma-s} apresenta o reporte da estratégia de busca segundo os 16 itens da extensão PRISMA-S~\cite{rethlefsen2021prismas} ao PRISMA 2020~\cite{page2021prisma}.
Os itens cobrem nomeação das fontes, plataformas, limites, métodos de busca complementares, gerenciamento de referências e atualizações.

\begin{longtable}{p{0.5cm}p{4cm}p{9cm}}
\caption{Reporte PRISMA-S (16 itens) da estratégia de busca.}
\label{tab:prisma-s} \\
\toprule
\textbf{\#} & \textbf{Item} & \textbf{Descrição} \\
\midrule
\endfirsthead
\toprule
\textbf{\#} & \textbf{Item} & \textbf{Descrição} \\
\midrule
\endhead
\bottomrule
\endfoot
1 & Nomes das bases & Scopus e arXiv. \\
2 & Plataforma/interface & Scopus via API Elsevier (\texttt{scopus-mcp}); arXiv via API REST com paginação completa (\texttt{fetch\_arxiv.py}). \\
3 & Registros de estudos & Nenhum registro de estudos (por exemplo, PROSPERO) consultado. \\
4 & Recursos online & Nenhum site organizacional consultado separadamente. \\
5 & Rastreamento de citações & Forward e backward via Semantic Scholar a partir de sete artigos-semente; ResearchRabbit como alternativa manual. \\
6 & Contato com autores & Nenhum autor contactado. \\
7 & Outros métodos & Nenhum além dos descritos nos itens 1, 2 e 5. \\
8 & Estratégias completas & Strings exatas por base documentadas no protocolo de busca, com versionamento em \textit{git}. \\
9 & Limites e restrições & Scopus: \texttt{PUBYEAR > 2022}; arXiv: categorias \texttt{cs.SE}, \texttt{cs.AI}, \texttt{cs.CL} e filtro pós-busca \texttt{--year-from 2023}; idiomas inglês e português. \\
10 & Desenvolvimento da estratégia & Mapeamento de três blocos conceituais, construção das strings com OR/AND, validação por sete artigos-semente e dois ciclos de refinamento iterativo documentados no protocolo versionado. \\
11 & Revisão por pares & Estratégia revisada pelo orientador antes da execução final. \\
12 & Gerenciamento de registros & Zotero com Better BibTeX para gerenciamento; arquivos \texttt{.bib} versionados em \textit{git}. \\
13 & Total de registros & 1.007 dos bancos (Scopus: 39; arXiv: 968) e 186 únicos por rastreamento de citações; total de 1.193 antes da deduplicação. \\
14 & Desduplicação & Elicit Pro deduplica automaticamente na importação; remoção manual subsequente totalizando 83 duplicatas. \\
15 & Ferramentas de automação & Scopus via MCP, arXiv via script Python, triagem assistida por Elicit Pro, rastreamento via Semantic Scholar MCP. \\
16 & Atualizações & Busca a ser re-executada antes da submissão se decorridos mais de três meses desde a execução inicial. \\
\end{longtable}


% =====================================================================
% Metadados do manuscrito (PRISMA 2020 itens 24-27)
% =====================================================================
\section*{Registro e protocolo}
Esta revisão sistemática não foi registrada em bases como PROSPERO,
que concentra revisões da área biomédica.
O protocolo foi desenvolvido antes da execução das buscas e está
disponível no pacote público de replicação (DOI:
10.5281/zenodo.19463131), junto das emendas registradas durante a
condução.

\section*{Financiamento}
Esta pesquisa não recebeu financiamento específico de agências de
fomento públicas, privadas ou sem fins lucrativos.

\section*{Conflitos de interesse}
Os autores declaram não haver conflitos de interesse.

\section*{Disponibilidade de dados e materiais}
O protocolo, as planilhas de triagem (título/resumo e texto
completo) com as decisões e razões de exclusão, os formulários de
extração estruturada dos 28 estudos incluídos e os artefatos
intermediários da síntese estão publicados como pacote de replicação
no Zenodo sob licença CC-BY-4.0 (DOI:
\href{https://doi.org/10.5281/zenodo.19463131}{10.5281/zenodo.19463131}),
com espelho em
\url{https://github.com/alexandre-zatti/slr-code-agent-memory}.

\bibliography{referencias/estudos-incluidos}

\end{document}
